\specialchapter{Introduction}

Spectral theory aims to study the properties of certain normed algebras, the fundamental examples being the algebra of complex-valued continuous functions on a topological space, the algebra of endomorphisms of a Banach space, and the convolution algebra of functions integrable with respect to a Haar measure on a locally compact group. The scope of this study is not necessarily apparent from this simple indication, but to grasp its importance one need only observe that two contemporary mathematicians, unless they belong to the same school, will often have as common knowledge, besides linear algebra and the elementary part of differential and integral calculus, only certain aspects of Fourier analysis.

In the first two chapters of this Book, we develop the most fundamental ideas of spectral theories, and consider in greater detail the third example mentioned above, namely, the Fourier transform, in the setting of a locally compact abelian group.

The following chapters will first continue the general theory of compact linear maps, then that of operators on Hilbert spaces, this case including linear maps defined only on a subspace (generally dense) of such a space. As is well known, this theory is the basis of the mathematical formalism of Quantum Mechanics. Finally, we will present the elementary theory of unitary representations of topological groups, particularly compact ones, which we have already had occasion to use in Chapter IX of the book Lie Groups and Lie Algebras.

The titles of the first two chapters are the same as those of the first edition of this Book, published in 1967. They have been revised, corrected, and expanded.

For ease of reading, we indicate below the most important changes found in this edition, and in particular the additions.

For Chapter I:

- In no. 3 of §3, we introduce the notion of proper partial map, which allows the Gelfand correspondence between locally compact spaces and commutative C*-algebras to be stated in a functorial manner.

- In §4, we added nos. 13 and 14, which deal with holomorphic functional calculus in real or complex normable algebras, and in algebras without a unit element.

- In §6, the presentation of the continuous functional calculus has been reworked (no. 6). Nos. 9 to 11 are new; they consider the notion of a positive element in a general C*-algebra, approximate units of such algebras, and the quotient of a C*-algebra by a closed two-sided ideal.

- §7, concerning the most elementary spectral properties of endomorphisms of Banach spaces, is essentially new; numbers 7 and 8 partially take up certain numbers of the old paragraph 6. It is this new paragraph that will serve as the basis for the more in-depth study of endomorphisms of Hilbert spaces in Chapter IV.

For Chapter II:

- §1 on the Fourier transform has been revised in detail, for example no. 6, which deals with the functoriality of Pontryagin duality. In the new no. 9, we make more explicit the statements most important for the groups \(\mathbf{R}^n\) and \((\mathbf{R}/\mathbf{Z})^n\). Fourier theory on \(\mathbf{R}^n\), whose most natural framework lies in the theory of tempered distributions, is stated here through its essential properties, without giving the proofs, which will appear in Chapter IV.

- We have finally added a Fourier theory summary sheet that brings together the definitions and properties of the main function spaces involved in this theory, along with the essential formulas.

Many of the exercises are new. In particular, we note the considerable increase in the number of exercises in §1 of Chapter II, reflecting the mathematical importance of Fourier theory.

\chapter{Normed algebras}

Unless stated otherwise, the algebras considered in this chapter are assumed to be associative.

\section{SPECTRA AND CHARACTERS}

In this paragraph, the letter K denotes a commutative field. If E and F are K-vector spaces, we denote by \(E \otimes F = E \otimes_{K} F\) .

\subsection{Unital algebras}

A unital algebra over K is a pair \((A,e)\) where A is an algebra over K with identity element, and e is the identity element of A. Since e is uniquely determined by A, we will sometimes say, by abuse of language, that A is a unital algebra. If \((A,e)\) and \((A',e')\) are unital algebras, a unital morphism of \((A,e)\) into \((A',e')\) is a morphism \(\varphi\) from A to \(A'\) such that \(\varphi(e)=e'\). A unital subalgebra of \((A,e)\) is a pair \((A',e)\) , where \(A'\) is a subalgebra of A containing e.

We will often denote the identity element by 1.

Lemma 1. --- Let A be an algebra. For every idempotent j of A, the subspace jAj of A is the set of \(x \in A\) such that xj = jx = x. This is a subalgebra of A which admits the unit element j.

The proof is elementary.

\subsection{Spectrum of an element in a unital algebra}

DEFINITION 1. --- Let A be a unital algebra over K and e its identity element. For every \(x \in \mathrm{A}\) we call the spectrum of \(x\) relative to A the set of \(\lambda \in \mathrm{K}\) such that \(\lambda e - x\) is not invertible.

The spectrum of x will be denoted \(\mathrm{Sp}_{\mathrm{A}}(x)\) , or \(\mathrm{Sp}(x)\) if no confusion results. The complement of \(\mathrm{Sp}_{\mathrm{A}}(x)\) in K is called the resolvent set of x.

Remarks. --- 1) If A = \{0\}, then we have Sp(0) = ∅.

\begin{enumerate}
\def\labelenumi{\arabic{enumi})}
\setcounter{enumi}{1}
\item
  If \(A \neq \{0\}\) , we have \(\mathrm{Sp}(\lambda e) = \{\lambda\}\) for all \(\lambda \in K\) .
\item
  For \(x \in \mathrm{A}\) to be invertible, it is necessary and sufficient that \(0 \notin \operatorname{Sp}(x)\) .
\item
  Let \(R \in K(X)\) be a rational fraction and let \(x \in A\) be an element that is substitutable in R, that is (A, IV, p.~20) there exist P and \(Q \in K[X]\) such that \(R = P/Q\) and \(Q(x)\) is invertible; one can then form the element \(R(x) = P(x) \cdot Q(x)^{-1} = Q(x)^{-1} \cdot P(x)\) of A; it does not depend on the choices of P and Q. We have \(0 \notin Q(\mathrm{Sp}(x))\) , so that every element of \(\mathrm{Sp}(x)\) is substitutable in R.
\end{enumerate}

We have \(\mathrm{R}(\mathrm{Sp}(x)) \subset \mathrm{Sp}(\mathrm{R}(x))\). Indeed, let \(\lambda \in \mathrm{Sp}(x)\) ; there exists a polynomial \(P_{1}\) such that \(\mathrm{R}(\lambda) - \mathrm{R}(\mathrm{X}) = (\lambda - \mathrm{X})\mathrm{P}_{1}(\mathrm{X}) / \mathrm{Q}(\mathrm{X})\) ; then, \(\mathrm{R}(\lambda)e - \mathrm{R}(x) = (\lambda e - x)(\mathrm{P}_{1}(x)/\mathrm{Q}(x))\) , so that \(\mathrm{R}(\lambda) - \mathrm{R}(x)\) is not invertible, hence \(\mathrm{R}(\lambda) \in \mathrm{Sp}(\mathrm{R}(x))\) .

Conversely, suppose that the field K is algebraically closed. Suppose first that R is not constant and let us prove that one has \(\mathrm{Sp}(\mathrm{R}(x)) = \mathrm{R}(\mathrm{Sp}(x))\) . Let \(\mu \in \mathrm{Sp}(\mathrm{R}(x))\) . Since R is not constant, P - \(\mu Q\) is not the zero polynomial; let \(\mu Q - P = \alpha \prod (\lambda_i - X)\) be a decomposition into factors of degree 1, so that \(\mu e - R(x) = \alpha \prod (\lambda_i e - x) Q(x)^{-1}\) . Since \(\mu e - R(x)\) is not invertible, there exists i such that \(\lambda_i e - x\) is not invertible, hence \(\lambda_i \in \mathrm{Sp}(x)\) , then \(\mathrm{R}(\lambda_i) = \mu \in \mathrm{R}(\mathrm{Sp}(x))\) .

When R is constant, the equality \(\mathrm{Sp}(\mathrm{R}(x)) = \mathrm{R}(\mathrm{Sp}(x))\) also holds, provided that \(\mathrm{Sp}(x)\) is nonempty.

\begin{enumerate}
\def\labelenumi{\arabic{enumi})}
\setcounter{enumi}{4}
\item
  Suppose that the algebra A is nonzero. Let \(x \in A\) be a nilpotent element. Let n be an integer such that \(x^{n} = 0\). The spectrum of \(x^{n}\) is reduced to 0, hence so is the spectrum of x by remark 4.
\item
  Let A and B be unital algebras over K and \(\varphi\colon\mathrm{A}\to\mathrm{B}\) a unital morphism. For every \(x\in\mathrm{A}\) , one has \(\mathrm{Sp}_{\mathrm{B}}(\varphi(x))\subset\mathrm{Sp}_{\mathrm{A}}(x)\) .
\item
  Let A be a unital algebra, R its radical (A, VIII, p.~150, def. 2) and let \(\varphi\) be the canonical morphism from A onto \(B = A/R\) . If \(x \in A\) , one has \(\mathrm{Sp}_{\mathrm{B}}(\varphi(x)) = \mathrm{Sp}_{\mathrm{A}}(x)\) . Indeed, it suffices to prove that if \(\varphi(x)\) is invertible in B, then \(x\) is invertible in A. Now, if \(y \in A\) is such that \(\varphi(x)\varphi(y) = \varphi(y)\varphi(x) = \varphi(e)\) , one has \(xy \in e + R\) , \(yx \in e + R\) , hence \(xy\) and \(yx\) are invertible (A, VIII, p.~151, th. 1) and consequently \(x\) is invertible. In particular, if \(x \in R\) , one has \(\mathrm{Sp}_{\mathrm{A}}(x) = \{0\}\) if \(A \neq \{0\}\) .
\item
  Let \((\mathrm{B}_{i})_{i\in\mathrm{I}}\) be a family of unital algebras, with \(\mathrm{B}_{i}=(\mathrm{A}_{i},e_{i})\) for \(i\in I\) . Set \(A=\prod_{i}A_{i}, e=(e_{i})_{i\in\mathrm{I}}\) . Then \((\mathrm{A},e)\) is a unital algebra called the product of the \(B_{i}\) . If \(x=(x_{i})_{i\in\mathrm{I}}\in\mathrm{A}\) , one has \(\mathrm{Sp}_{\mathrm{A}}(x)=\bigcup_{i}\mathrm{Sp}_{\mathrm{A}_{i}}(x_{i})\) .
\end{enumerate}

Examples. --- 1) Let X be a set and let A = K \(^{X}\) be the algebra of functions with values in K defined on X. The spectrum of an element f of A is the set of values of f.

\begin{enumerate}
\def\labelenumi{\arabic{enumi})}
\setcounter{enumi}{1}
\tightlist
\item
  Let A be a unital algebra of finite rank over K. For \(x \in A\) to be invertible, it is necessary and sufficient that the linear map \(y \mapsto xy\) from A to A have nonzero determinant. It follows that the spectrum of x is the set of roots of the characteristic polynomial of x (A, III, p.~110). If A is the algebra of endomorphisms of a vector space V of finite dimension over K, the spectrum of x is therefore the set of eigenvalues of x. This is not always the case when V is of infinite dimension (cf.~I, p.~153, exercise 2).
\end{enumerate}

\subsection{Resolvent}

DEFINITION 2. --- Let A be a unital algebra over K and \(x \in \mathbf{A}\) . For every \(\lambda \in \mathrm{K} - \mathrm{Sp}(x)\) , set

\[
\mathrm{R} (x, \lambda) = (\lambda e - x) ^ {- 1}.
\]

The map from \(K-\mathrm{Sp}(x)\) to A given by \(\lambda\mapsto\mathrm{R}(x,\lambda)\) is called the resolvent of \(x\) .

For \(x\) fixed, the values of \(R(x, \lambda)\) are pairwise permutable. If \(\lambda, \mu \in K\) , we have:

\[
(\lambda - \mu) e = (\lambda e - x) - (\mu e - x)
\]

hence, if \(\lambda, \mu \in \mathrm{K} - \mathrm{Sp}(x)\) , one has the relation

\[
(\lambda - \mu) \mathrm{R} (x, \lambda) \mathrm{R} (x, \mu) = \mathrm{R} (x, \mu) - \mathrm{R} (x, \lambda).\tag{1}
\]

If \(x,y\in \mathrm{A}\) and \(\lambda \in \mathrm{K}\) , one has:

\[
y - x = (\lambda e - x) - (\lambda e - y)
\]

hence, if \(\lambda\in\mathrm{K}-(\mathrm{Sp}(x)\cup\mathrm{Sp}(y))\) , one has the relation

\[
\mathrm{R} (y, \lambda) (y - x) \mathrm{R} (x, \lambda) = \mathrm{R} (y, \lambda) - \mathrm{R} (x, \lambda).\tag{2}
\]

\subsection{Spectrum of an element in an algebra}

Let A be an algebra over K. Recall (A, III, p.~4) that one defines on the vector space \(\widetilde{\mathbf{A}} = \mathbf{K} \times \mathbf{A}\) an algebra structure such that:

\[
(\lambda , a) (\mu , b) = (\lambda \mu , \lambda b + \mu a + a b).
\]

Let \(e = (1, 0)\) . Then \((\tilde{\mathrm{A}}, e)\) is the unital algebra obtained from A by adjoining a unit element. The algebra A is identified with the two-sided ideal \(\{0\} \times A\) of \(\tilde{A}\) ; the algebra A is commutative if and only if \(\tilde{A}\) is.

If A' is a second algebra over K, let \((\tilde{\mathrm{A}}', e')\) be the unital algebra obtained from A' by adjoining an identity element, and let \(\varphi\) be a morphism from A to A'; there exists one and only one unital morphism from \((\tilde{\mathrm{A}}, e)\) into \((\tilde{\mathrm{A}}', e')\) which extends \(\varphi\) .

Let A be an algebra over K and \(x \in A\) . We call the spectrum of x relative to A the spectrum of x relative to the unital algebra \(\widetilde{A}\) deduced from A by adjoining an identity element. This set will be denoted \(\mathrm{Sp}_{\mathrm{A}}^{\prime}(x)\) , or \(\mathrm{Sp}^{\prime}(x)\) if no confusion arises. We have \(0 \in \mathrm{Sp}_{\mathrm{A}}^{\prime}(x)\) for all \(x \in A\) .

If \(\varphi\) is a morphism from A into an algebra B, we have \(\mathrm{Sp}_{\mathrm{B}}^{\prime}(\varphi(x)) \subset \mathrm{Sp}_{\mathrm{A}}^{\prime}(x)\) .

Remarks. --- 1) Let (A,1) be a unital algebra. If \(x \in A\) , one has

\[
\operatorname{Sp} _ {\mathrm{A}} ^ {\prime} (x) = \operatorname{Sp} _ {\mathrm{A}} (x) \cup \{0 \}.
\]

We verify in fact that \((e-1) \cdot A = A \cdot (e-1) = 0\) , so that \(\widetilde{A}\) is the unital algebra product of A and \(K(e-1)\) . Our assertion therefore follows from Remark 8 of I, p.~3.

\begin{enumerate}
\def\labelenumi{\arabic{enumi})}
\setcounter{enumi}{1}
\tightlist
\item
  It follows from Remark 1 that, if B is an algebra over K and if \(x \in B\) , we have:
\[
\mathrm{Sp} _ {\mathrm{B}} ^ {\prime} (x) = \mathrm{Sp} _ {\widetilde {\mathrm{B}}} (x) = \mathrm{Sp} _ {\widetilde {\mathrm{B}}} (x) \cup \{0 \} = \mathrm{Sp} _ {\widetilde {\mathrm{B}}} ^ {\prime} (x).
\]

\item
  If x belongs to the radical of A (A, VIII, p.~430, def. 3), then \(\mathrm{Sp}_{\mathrm{A}}^{\prime}(x)=\{0\}\). This follows from Remark 7 of I, p.~3.
\end{enumerate}

PROPOSITION 1. --- Let A be an algebra and \(x, y \in \mathrm{A}\) . We have \(\mathrm{Sp}'(xy) = \mathrm{Sp}'(yx)\) .

Passing to \(\widetilde{\mathbf{A}}\) we reduce to the case where A has a unit element \(e\) . It then suffices to prove that, if \(\lambda \neq 0\) is such that \(xy - \lambda e\) admits an inverse \(u\) , then \(yx - \lambda e\) is invertible. Set \(z = yux - e\) . Since \(xyu = \lambda u + e\) , one has

\[
\begin{array}{r l} (y x - \lambda e) z & = y (x y u) x - y x - \lambda y u x + \lambda e \\ & = y (\lambda u + e) x - y x - \lambda y u x + \lambda e = \lambda e \end{array}
\]

and likewise \(z(yx - \lambda e) = \lambda e\) . Since \(\lambda \neq 0\) , we see that \(yx - \lambda e\) is invertible.

If A is a unital algebra and if \(x, y \in A\) , the preceding proposition implies that \(\operatorname{Sp}(xy) \cup \{0\} = \operatorname{Sp}(yx) \cup \{0\}\) , but one may have \(\operatorname{Sp}(xy) \neq \operatorname{Sp}(yx)\) (cf.~I, p.~153, exerc. 3).

\subsection{Full subalgebras}

Let A be a unital algebra over K.

DEFINITION 3. --- A full subalgebra of A is a unital subalgebra B such that every element of B which is invertible in A is invertible in B.

In other words, B is a full subalgebra of A if and only if \(\mathrm{Sp}_{\mathrm{B}}(x)=\mathrm{Sp}_{\mathrm{A}}(x)\) for all \(x\in B\) .

The intersection of a family of full subalgebras of A is a full subalgebra of A.

Let M be a subset of A. The intersection of the full subalgebras of A containing M is the smallest full subalgebra of A containing M.

It is called the full subalgebra of A generated by M. The commutant M′ of M in A is a full subalgebra of A (for, if x is invertible in A and commutes with M, then \(x^{-1}\) commutes with M). Therefore, the bicommutant M'' of M is a full subalgebra of A that contains the full subalgebra of A generated by M.

If the elements of M are pairwise permutable, we have \(M \subset M'\) and \(M'' \subset M'''\) ; the algebra \(M''\) is therefore commutative, and the same is then true of the full subalgebra generated by M.

A maximal commutative subalgebra of A is a full subalgebra, because it is equal to its commutant.

Lemma 2. --- Let \((x_{\lambda})_{\lambda \in \Lambda}\) be a family of pairwise commuting elements of A. The full subalgebra generated by \((x_{\lambda})\) is the set of elements of the form \(\mathrm{R}((x_{\lambda}))\) , where \(\mathrm{R} \in \mathrm{K}((\mathrm{X}_{\lambda}))\) ranges over the set of rational fractions in which the family \((x_{\lambda})\) is substitutable.

Let B be the full subalgebra of A generated by the family \((x_{\lambda})\) , and denote by \(B_{1}\) the set of elements of the form \(\mathrm{R}((x_{\lambda}))\) , where \(\mathrm{R}((\mathrm{X}_{\lambda}))\) is a rational fraction in which \((x_{\lambda})\) is substitutable. Explicitly, \(B_{1}\) is the set of elements of A of the form \(\mathrm{P}((x_{\lambda}))\mathrm{Q}((x_{\lambda}))^{-1}\) , where \(\mathrm{P},\mathrm{Q}\in\mathrm{K}[(\mathrm{X}_{\lambda})]\) and \(\mathrm{Q}((x_{\lambda}))\) is invertible in A. The set \(B_{1}\) is a unital subalgebra of A containing the family \((x_{\lambda})\) . It is a full subalgebra: if \(\mathrm{P}((x_{\lambda}))\mathrm{Q}((x_{\lambda}))^{-1}\) is invertible in A, then \(\mathrm{P}((x_{\lambda}))\) is invertible in A and the inverse \(\mathrm{Q}((x_{\lambda}))\mathrm{P}((x_{\lambda}))^{-1}\) of \(\mathrm{P}((x_{\lambda}))\mathrm{Q}((x_{\lambda}))^{-1}\) belongs to \(B_{1}\) . We therefore have \(B\subset B_{1}\) . On the other hand, if \(P,Q\in K[(X_{\lambda})]\) , and if \(\mathrm{Q}((x_{\lambda}))\) is invertible in A, then \(\mathrm{P}((x_{\lambda}))\in\mathrm{B}\) and \(\mathrm{Q}((x_{\lambda}))\in\mathrm{B}\) , hence \(\mathrm{Q}((x_{\lambda}))^{-1}\in\mathrm{B}\) and \(\mathrm{P}((x_{\lambda}))\mathrm{Q}((x_{\lambda}))^{-1}\in\mathrm{B}\) , hence \(B_{1}\subset B\) .

\subsection{Characters of a commutative unital algebra}

DEFINITION 4. --- Let A be a commutative unital algebra over K. A unital character is a unital morphism from A to K.

When no confusion can arise, one simply says character instead of unital character. The set of unital characters of A is denoted X(A). If A is the zero algebra, then X(A) is empty.

Let A and B be commutative unital algebras over K and h a unital morphism from A to B. The map \(\chi \mapsto \chi \circ h\) of \(\mathsf{X}(\mathsf{B})\) into \(\mathsf{X}(\mathsf{A})\) is denoted \(\mathsf{X}(h)\) . If k is a morphism from B to a commutative unital algebra, one has \(\mathsf{X}(k \circ h) = \mathsf{X}(h) \circ \mathsf{X}(k)\) . The map \(\mathsf{X}(\mathrm{Id}_{\mathrm{A}})\) is the identity map of \(\mathsf{X}(\mathrm{A})\) .

If h is surjective, \(\mathsf{X}(h)\) is a bijection from \(\mathsf{X}(\mathsf{B})\) onto the set of characters of A that vanish on the kernel of h.

Let \((\mathrm{A}_{1},e_{1}),\ldots,(\mathrm{A}_{n},e_{n})\) be commutative unital algebras over K and let A be the unital algebra \(A_{1}\times\cdots\times A_{n}\) , with identity element \((e_{1},\ldots,e_{n})\) . For each i, identify \(A_{i}\) with an ideal of A and let \(\pi_{i}\) be the canonical map from A onto \(A_{i}\) . Then \(\mathsf{X}(\pi_{i})\) is a bijection from \(\mathsf{X}(\mathrm{A}_{i})\) onto the set \(X_{i}\) of characters of A vanishing on \(\prod_{j\neq i}A_{j}\) . The sets \(X_{i}\) are pairwise disjoint. On the other hand, let \(\chi\in\mathsf{X}(\mathrm{A})\) . Since \(1=\sum\chi(e_{i})\) , there exists i such that \(\chi(e_{i})\neq0\) . For every \(j\neq i\) and every \(y\in A_{j}\) , one has \(\chi(e_{i})\chi(y)=\chi(e_{i}y)=\chi(0)=0\) , hence \(\chi(\mathrm{A}_{j})=0\). Thus, \(\chi\) vanishes on \(\prod_{j\neq i}A_{j}\) , so that \(\mathsf{X}(\mathrm{A})\) is the union of the \(X_{i}\) .

Let B be the unital algebra \(A_{1} \otimes \cdots \otimes A_{n}\) . Denote by \(h_{i}\) the canonical morphism \(A_{i} \to B\) . Then

\[
\chi \mapsto (\chi \circ h _ {1}, \dots , \chi \circ h _ {n})
\]

is a map from \(\mathsf{X}(\mathsf{B})\) into \(\mathsf{X}(\mathsf{A}_1)\times \dots \times \mathsf{X}(\mathsf{A}_n)\) , and

\[
\left(\chi_ {1}, \dots , \chi_ {n}\right) \mapsto \chi_ {1} \otimes \dots \otimes \chi_ {n}
\]

is a map from \(\mathsf{X}(\mathsf{A}_{1}) \times \cdots \times \mathsf{X}(\mathsf{A}_{n})\) into \(\mathsf{X}(\mathsf{B})\) . One verifies that these maps are inverse bijections of one another, by which one identifies \(\mathsf{X}(\mathsf{B})\) with \(\mathsf{X}(\mathsf{A}_{1}) \times \cdots \times \mathsf{X}(\mathsf{A}_{n})\) .

Let A be a commutative unital algebra over K. Let Y be the set of ideals of codimension 1 of A. For every \(\chi \in \mathsf{X}(\mathsf{A})\) , one has \(\operatorname{Ker}(\chi) \in \mathsf{Y}\) . The map \(\chi \mapsto \operatorname{Ker}(\chi)\) is a bijection from \(\mathsf{X}(\mathsf{A})\) onto Y. Indeed, if \(I \in Y\) , there exists a unique isomorphism of the unital K-algebra A/I onto K and the composite morphism

\[
\mathrm{A} \longrightarrow \mathrm{A} / \mathrm{I} \longrightarrow \mathrm{K}
\]

is the unique character of A with kernel I.

DEFINITION 5. --- Let A be a commutative unital algebra over K. For every \(x \in \mathrm{A}\) , we denote by \(\mathcal{G}_{\mathrm{A}}(x)\) , or simply \(\mathcal{G}(x)\) , the map \(\chi \mapsto \chi(x)\) from \(\mathrm{X}(\mathrm{A})\) to K. It is called the Gelfand transform of \(x\) .

The map \(\mathcal{G}\) is a unital morphism from A into the unital algebra \(\mathrm{K}^{\mathsf{X}(A)}\) of maps from \(\mathsf{X}(A)\) to K. It is called the Gelfand transform of A.

If \(x \in \mathrm{A}\) , the image of the Gelfand transform \(\mathcal{G}_{\mathrm{A}}(x)\) of \(x\) is contained in \(\mathrm{Sp}_{\mathrm{A}}(x)\) . Indeed, let \(\chi \in \mathsf{X}(\mathrm{A})\) ; since \(\chi(x - \chi(x)e) = 0\) , the element \(x - \chi(x)e\) is not invertible.

Let B be a unital commutative algebra over K and h a unital morphism from A to B; then \(\mathsf{X}(h)\colon\mathsf{X}(\mathsf{B})\to\mathsf{X}(\mathsf{A})\) defines a unital morphism \(h_{*}\colon\mathsf{K}^{\mathsf{X}(\mathsf{A})}\to\mathsf{K}^{\mathsf{X}(\mathsf{B})}\) , and the diagram:

\[
\begin{array}{l} \text {A} \xrightarrow {\mathcal {G} _ {\text {A}}} \text {K} ^ {\mathsf {X} (\text {A})} \\ \Big \downarrow h \qquad \qquad \qquad \Big \downarrow h _ {*} \\ \text {B} \xrightarrow {\mathcal {G} _ {\text {B}}} \text {K} ^ {\mathsf {X} (\text {B})} \end{array}\tag{3}
\]

is commutative. Indeed, for every \(x \in A\) and every \(\chi \in \mathsf{X}(\mathsf{B})\) , we have:

\[
\begin{array}{r l} & {\mathcal {G} _ {\mathrm{B}} (h (x)) (\chi) = \chi (h (x)) = (\chi \circ h) (x)} \\ & {\qquad = (\mathsf {X} (h) (\chi)) (x)} \\ & {\qquad = \mathcal {G} _ {\mathrm{A}} (x) (\mathsf {X} (h) (\chi))} \\ & {\qquad = h _ {*} (\mathcal {G} _ {\mathrm{A}} (x)) (\chi).} \end{array}\tag{4}
\]

Suppose now that K is a topological field. We then equip \(\mathsf{X}(\mathsf{A})\) with the topology of pointwise convergence on A (cf.~EVT, III, p.~14, example 1), and the topological space \(\mathsf{X}(\mathsf{A})\) is called the character space of A. The topology of \(\mathsf{X}(\mathsf{A})\) is therefore the coarsest for which the functions \(\mathcal{G}_{\mathrm{A}}(x)\) for \(x \in A\) are continuous, and the map \(\chi \mapsto (\chi(a))_{a \in \mathrm{A}}\) identifies the space \(\mathsf{X}(\mathsf{A})\) with a subset of \(K^{A}\) .

When K = R or C, this topology is the topology induced on \(\mathsf{X}(\mathsf{A}) \subset \mathsf{A}^{*}\) by the weak topology \(\sigma(\mathsf{A}^{*}, \mathsf{A})\) on \(A^{*}\) (EVT, II, p.~45, def. 2); as such, we will also say that it is the weak topology on \(\mathsf{X}(\mathsf{A})\) .

If \(h\) is a unital morphism from A to B, the map \(\mathsf{X}(h)\colon \mathsf{X}(\mathsf{B})\to \mathsf{X}(\mathsf{A})\) is continuous. If \(h\) is surjective, the image of \(\mathsf{X}(h)\), namely, the set of characters of A vanishing on the kernel of \(h\), is closed in \(\mathsf{X}(\mathsf{A})\) ; on the other hand, the topology on \(\mathsf{X}(h)(\mathsf{X}(\mathsf{B}))\) deduced from that of \(\mathsf{X}(\mathsf{B})\) by the bijection \(\mathsf{X}(h)\) is the topology of pointwise convergence in A, that is, the topology induced by that of \(\mathsf{X}(\mathsf{A})\) . In other words, \(\mathsf{X}(h)\) is a homeomorphism from \(\mathsf{X}(\mathsf{B})\) onto a closed subset of \(\mathsf{X}(\mathsf{A})\) .

If \(A_{1},\ldots,A_{n}\) are unital commutative algebras over K, the space \(\mathsf{X}(\mathsf{A}_{1}\times\cdots\times\mathsf{A}_{n})\) is thus identified with the topological sum space of \(\mathsf{X}(\mathsf{A}_{1}),\ldots,\mathsf{X}(\mathsf{A}_{n})\) . Likewise, \(\mathsf{X}(\mathsf{A}_{1}\otimes\cdots\otimes\mathsf{A}_{n})\) is identified with the product topological space \(\mathsf{X}(\mathsf{A}_{1})\times\cdots\times\mathsf{X}(\mathsf{A}_{n})\) .

\subsection{Case of algebras without a unit element}

DEFINITION 6. --- Let A be a commutative algebra over K. A character of A is an algebra homomorphism from A to K.

The set of characters of A will be denoted by X'(A).

The zero map is an algebra homomorphism. If A has a unit element, \(e\) , a nonzero algebra morphism from A into K is unital, that is, it is a unital character of A in the sense of Definition 4: indeed, for \(\chi \in \mathsf{X}'(\mathsf{A})\), \(\chi\) is nonzero if and only if \(\chi(e) = 1\) .

We set \(\mathsf{X}(\mathsf{A})=\mathsf{X}^{\prime}(\mathsf{A})-\{0\}\) ; by the preceding, the notation is compatible with that introduced when A is unital.

If \(h\colon A\to B\) is a morphism of commutative algebras, the map \(\chi\mapsto\chi\circ h\) is a map \(\mathsf{X}^{\prime}(h)\colon\mathsf{X}^{\prime}(\mathrm{B})\to\mathsf{X}^{\prime}(\mathrm{A})\) . It maps 0 to 0. If \(k\colon B\to C\) is a morphism of commutative algebras, then one has \(\mathsf{X}^{\prime}(k\circ h)=\mathsf{X}^{\prime}(h)\circ\mathsf{X}^{\prime}(k)\) . If \(h\) is surjective, \(\mathsf{X}^{\prime}(h)\) is a bijection from \(\mathsf{X}^{\prime}(\mathrm{B})\) onto the set of characters of A vanishing on the kernel of \(h\) . Let \(\mathrm{A}_{1},\ldots,\mathrm{A}_{n}\) be commutative algebras, \(\mathrm{A}=\mathrm{A}_{1}\times\cdots\times\mathrm{A}_{n}\) and \(\pi_{i}\colon\mathrm{A}\to\mathrm{A}_{i}\) the canonical morphism; then \(\mathsf{X}^{\prime}(\pi_{i})\) is a bijection from \(\mathsf{X}^{\prime}(\mathrm{A}_{i})\) onto a subset \(\mathsf{X}_{i}^{\prime}\) of \(\mathsf{X}^{\prime}(\mathrm{A})\) , namely the set of characters of A vanishing on \(\prod_{j\neq i}\mathrm{A}_{j}\) ; we see as in no. 6 that \(\mathsf{X}^{\prime}(\mathrm{A})\) is the union

of \(X_{i}^{\prime}\) ; on the other hand, \(X_{i}^{\prime} \cap X_{j}^{\prime} = \{0\}\) for \(i \neq j\) ; in particular the \(X_{i}^{\prime} - \{0\}\) form a partition of \(X^{\prime}(A) - \{0\} = X(A)\) .

For every \(x \in A\) , let \(\mathcal{G}_{\mathrm{A}}^{\prime}(x)\) , or simply \(\mathcal{G}^{\prime}(x)\) , be the map \(\chi \mapsto \chi(x)\) from \(\mathsf{X}^{\prime}(\mathrm{A})\) to K. The map \(\mathcal{G}^{\prime}\) is a morphism from A into the algebra \(A_{1}\) of the maps \(\mathsf{X}^{\prime}(\mathrm{A}) \to \mathrm{K}\) vanishing at 0. Let B be a commutative algebra, let \(B_{1}\) be the algebra of maps \(\mathsf{X}^{\prime}(\mathrm{B}) \to \mathrm{K}\) vanishing at 0, and let h be a morphism from A to B; then \(\mathsf{X}^{\prime}(h)\) defines a morphism \(h_{1}: A_{1} \to B_{1}\) , and we have \(h_{1} \circ \mathcal{G}_{A}^{\prime} = \mathcal{G}_{B}^{\prime} \circ h\) . One denotes by \(\mathcal{G}_{\mathrm{A}}(x)\) , or simply \(\mathcal{G}(x)\) the restriction of \(\mathcal{G}_{\mathrm{A}}^{\prime}(x)\) onto \(\mathrm{X}(\mathrm{A})\) and it is called the Gelfand transform of \(x\) .

Let \(\widetilde{\mathsf{A}}\) be the unital algebra deduced from A by adjoining an identity element. By restriction, every character of \(\widetilde{\mathsf{A}}\) defines a character of A; conversely, every character of A extends uniquely to a character of \(\widetilde{\mathsf{A}}\) . This defines a canonical bijection from \(\mathsf{X}'(\mathsf{A})\) onto \(\mathsf{X}(\widetilde{\mathsf{A}})\) , by which these two sets are identified. The character 0 of A is identified with the unique character of \(\widetilde{\mathsf{A}}\) with kernel A.

If \(x\in \mathrm{A}\) and \(\chi \in \mathsf{X}'(\mathrm{A})\) , one has \(\chi (x)\in \mathrm{Sp}_{\widetilde{\mathrm{A}}}^{\sim}(x)\) , hence \(\chi (x)\in \mathrm{Sp}_{\mathrm{A}}^{\prime}(x)\) .

Lemma 3. --- The map \(\chi \mapsto \mathrm{Ker}(\chi)\) is a bijection from \(\mathsf{X}(\mathsf{A})\) onto the set of regular ideals of codimension 1 of A.

Recall (A, VIII, p.~426, def. 1) that an ideal I of A is said to be regular if the quotient algebra A/I admits a unit element.

Let us prove the lemma. On the one hand, \(\mathsf{X}(\mathsf{A})\) is identified with the set of characters of \(\widetilde{A}\) nonzero on A. On the other hand, by A, VIII, p.~428, prop. 4, the map \(I \mapsto A \cap I\) is a bijection from the set of maximal ideals of \(\widetilde{A}\) distinct from A onto the set of regular maximal ideals of A. The lemma then follows from the results of no. 6.

Suppose now that K is a topological field. We then equip \(\mathsf{X}^{\prime}(\mathsf{A})\) with the topology of pointwise convergence on A; the notation \(\mathsf{X}^{\prime}(\mathsf{A})\) will henceforth denote the topological space thus obtained. When K = R or C, we shall also call it the weak topology. For every \(x \in A\) , the function \(\mathcal{G}_{\mathrm{A}}^{\prime}(x)\) on \(\mathsf{X}^{\prime}(\mathsf{A})\) is continuous.

If h is a morphism from A to B, the map \(\mathsf{X}'(h)\colon\mathsf{X}'(\mathsf{B})\to\mathsf{X}'(\mathsf{A})\) is continuous. If h is surjective, \(\mathsf{X}'(h)\) is a homeomorphism from \(\mathsf{X}'(\mathsf{B})\) onto its image and this image is closed in \(\mathsf{X}'(\mathsf{A})\) .

Let \(A = A_{1} \times \cdots \times A_{n}\) ; with the same notations as above, \(\mathsf{X}'(\pi_{i})\) is a homeomorphism from \(\mathsf{X}'(\mathrm{A}_{i})\) onto \(X_{i}'\) and \(X_{i}'\) is closed in \(\mathsf{X}'(\mathrm{A})\) . Therefore \(X_{i}' - \{0\}\) is open in \(\mathsf{X}'(\mathrm{A})\) ; the \(\mathsf{X}'(\pi_{i})\) define a continuous map from the sum space S of the \(\mathsf{X}'(\mathrm{A}_{i})\) onto \(\mathsf{X}'(\mathrm{A})\) , and one verifies immediately that the image of a union of neighborhoods of the points \(0 \in \mathsf{X}'(\mathrm{A}_{1}), \ldots, 0 \in \mathsf{X}'(\mathrm{A}_{n})\) is a neighborhood of \(0 \in \mathsf{X}'(\mathrm{A})\) ; from all this it follows that \(\mathsf{X}'(\mathrm{A})\) is canonically identified with a quotient space of S. In particular, the space \(\mathsf{X}(\mathrm{A})\) is identified with the sum space of the \(\mathsf{X}(\mathrm{A}_{i})\) .

The canonical bijection from \(\mathsf{X}^{\prime}(\mathsf{A})\) onto \(\mathsf{X}(\tilde{\mathsf{A}})\) is a homeomorphism. Let B be a unital algebra over K and \(B^{\prime}\) the underlying algebra; then the space \(\mathsf{X}(\mathsf{B})\) is identified with the subspace \(\mathsf{X}(\mathsf{B}^{\prime})\) of \(\mathsf{X}^{\prime}(\mathsf{B}^{\prime})\) .

\subsection{Primitive ideals}

Let A be an algebra over K and E a vector space over K. One calls a representation of A in E a morphism from A to the algebra \(\mathcal{L}(\mathrm{E})\) of endomorphisms of E. An injective representation is said to be faithful. Let \(\pi_1\) and \(\pi_2\) be representations of A in spaces \(\mathrm{E}_1\) , \(\mathrm{E}_2\) . A morphism of \(\pi_1\) into \(\pi_2\) is a K-linear map \(u: \mathrm{E}_1 \to \mathrm{E}_2\) such that \(u(\pi_1(a)x) = \pi_2(a)u(x)\) for all \(a \in \mathrm{A}\) and \(x \in \mathrm{E}_1\) . Representations are said to be equivalent if there exists a morphism of \(\pi_1\) into \(\pi_2\) which is an isomorphism of vector spaces. Its inverse is then a morphism of \(\pi_2\) into \(\pi_1\) . A representation \(\pi\) of A in E is said to be irreducible if \(\mathrm{E} \neq \{0\}\) and if the only vector subspaces of E stable under \(\pi(\mathrm{A})\) are \(\{0\}\) and E.

Example. --- The zero map from A to \(\mathcal{L}(\mathrm{E})\) is a representation, called the trivial representation, of A. It is irreducible if and only if E has dimension 1.

Lemma 4. --- Let π be an irreducible, non-trivial representation of A in E. For every non-zero element ξ of E, one has π(A)ξ = E.

The subspace \(\pi(A)\xi\) of \(E\) is stable under \(\pi(A)\) . Suppose it is zero. The non-zero subspace \(K\xi\) of \(E\) would then be stable under \(\pi(A)\) , and hence equal to \(E\) ; but this would imply that \(\pi\) is the zero representation. We therefore have \(\pi(A)\xi = E\) .

Let \(\pi\) be an irreducible, non-trivial representation of A in E. By this lemma, the annihilator R of \(\xi\) in A is a regular left ideal (A, VIII, p.~425, n° 1) of A, and the representation \(\pi\) is equivalent to the representation defined by the A-pseudomodule A/R. Since \(\pi\) is irreducible, the ideal R is a maximal regular left ideal.

DEFINITION 7. --- Let A be an algebra over K. One calls a primitive ideal of A the kernel of a non-trivial irreducible representation of A.

If A is commutative, the primitive ideals of A are the maximal regular ideals of A. Indeed, the non-trivial irreducible representations of A are, up to equivalence, the representations \(\pi_{R}\) defined by the A-pseudomodules A/R, where R is a maximal regular ideal of A. The kernel of \(\pi_{R}\) contains R. It is even equal to it since, by A, VIII, p.~426, prop. 2, the commutativity of A implies that A/R is a field. Hence \(\text{Ker}(\pi_{\text{R}})\) is maximal regular.

Lemma 5. --- Let π be an irreducible representation of A in a vector space E over K.

\begin{enumerate}
\def\labelenumi{\alph{enumi})}
\item
  Let I be a two-sided ideal of A. If \(\pi(\mathrm{I}) \neq \{0\}\) , then \(\pi|\mathrm{I}\) is irreducible;
\item
  Let \(I_{1}\) and \(I_{2}\) be two-sided ideals of A such that \(\pi(\mathrm{I}_{1}) \neq 0\) and \(\pi(\mathrm{I}_{2}) \neq 0\) . Then \(\pi(\mathrm{I}_{1}\mathrm{I}_{2}) \neq 0\) .
\end{enumerate}

The set of elements of E annihilated by \(\pi(\mathrm{I})\) is stable under \(\pi(\mathrm{A})\) and distinct from E, hence equal to 0. Therefore, if \(\xi\) is a non-zero element of E, one has \(\pi(\mathrm{I})\xi \neq 0\) ; since \(\pi(\mathrm{I})\xi\) is stable under \(\pi(\mathrm{A})\) , one has \(\pi(\mathrm{I})\xi = \mathrm{E}\) , which proves \(a\) ). On the other hand, what precedes proves that \(\pi(\mathrm{I}_2)\mathrm{E} = \mathrm{E}\) , \(\pi(\mathrm{I}_1)\pi(\mathrm{I}_2)\mathrm{E} = \mathrm{E}\) , hence \(\pi(\mathrm{I}_1\mathrm{I}_2) \neq 0\) , whence \(b\) ).

Lemma 6. --- Let I \(_{1}\) and I \(_{2}\) be two-sided ideals of A, and I a primitive ideal of A. If I contains I \(_{1}\) I \(_{2}\) (in particular, if I contains I \(_{1}\) ∩ I \(_{2}\) ), then I contains I \(_{1}\) or I \(_{2}\) .

Let \(\pi\) be an irreducible representation with kernel I. If I \(\not\supseteq\) I \(_1\) and I \(\not\supseteq\) I \(_2\) , Lemma 5, \(b\) ) proves that \(\pi(\text{I}_1\text{I}_2) \neq 0\) , whence I \(\not\supseteq\) I \(_1\) I \(_2\) .

Lemma 7. --- Suppose that A has a unit element. Let I be a maximal two-sided ideal of A. Then I is a primitive ideal.

There exists a maximal left ideal R of A containing I (A, I, p.~99, th. 1). Let \(\pi\) be the canonical representation of A in A/R, which is irreducible and nonzero. Since IA \(\subset\) R, the kernel I' of \(\pi\) contains I, hence I' = I and I is primitive.

Let J(A) be the set of primitive ideals of A. For any subset M of A, we denote by V(M) the set of primitive ideals of A containing M; if I is the two-sided ideal of A generated by M, then V(M) = V(I). If M is reduced to a single element x, we write V(x) instead of V(\{x\}).

The mapping M \(\mapsto\) V(M) is decreasing with respect to inclusion relations. We have:

\[
\begin{array}{r l} (5) & \mathrm{V} (\varnothing) = \mathrm{J} (\mathrm{A}), \qquad \mathrm{V} (\mathrm{A}) = \varnothing \end{array}
\]

\[
\begin{array}{r l} (6) & \mathrm{V} \left(\bigcup_ {i \in \mathrm{I}} \mathrm{M} _ {i}\right) = \mathrm{V} \left(\sum_ {i \in \mathrm{I}} \mathrm{M} _ {i}\right) = \bigcap_ {i \in \mathrm{I}} \mathrm{V} (\mathrm{M} _ {i}) \end{array}
\]

for any family \((\mathrm{M}_{i})_{i\in\mathrm{I}}\) of subsets of A. On the other hand, by Lemma 6,

\[
\begin{array}{r l} (7) & \mathrm{V} (\mathrm{I} _ {1} \cap \mathrm{I} _ {2}) = \mathrm{V} (\mathrm{I} _ {1} \mathrm{I} _ {2}) = \mathrm{V} (\mathrm{I} _ {1}) \cup \mathrm{V} (\mathrm{I} _ {2}) \end{array}
\]

for all two-sided ideals \(I_{1}\) , \(I_{2}\) of A. Formulas (5) through (7) show that the subsets V(M) of J(A) are the closed subsets of a topology called the Jacobson topology on J(A).

Let T be a subset of J(A), and let \(\Upsilon(T)\) be the intersection of the elements of T, so that \(\Upsilon(T)\) is a two-sided ideal of A. Then the closure of T in J(A) is the smallest closed subset of J(A) containing T, that is, V( \(\Upsilon(T)\) ). In particular, T is closed if and only if \(T = V(\Upsilon(T))\) .

PROPOSITION 2. --- Let I \(_{1}\) and I \(_{2}\) be distinct points of J(A). Then one of these two points is not adherent to the other.

Indeed, we have, for example, \(I_{1} \not\subset I_{2}\) . The set \(V(I_{1})\) of \(I \in J(A)\) such that \(I_{1} \subset I\) is closed in \(J(A)\) and it contains \(I_{1}\) but not \(I_{2}\) .

PROPOSITION 3. --- Let I ∈ J(A). For \{I\} to be closed in J(A), it is necessary and sufficient that I be a primitive maximal ideal.

Indeed, the closure of \{I\} consists of the primitive ideals of A containing I.

The relation

« \(\pi_{1},\pi_{2}\) are representations of A, which are isomorphic.

is an equivalence relation with respect to \(\pi_1\) and \(\pi_2\) . For any representation \(\pi\) of A, we denote by \(\operatorname{cl}(\pi)\) the equivalence class of \(\pi\) , which is therefore a representation of A isomorphic to \(\pi\) , such that two representations \(\pi_1\) and \(\pi_2\) are isomorphic if and only if \(\operatorname{cl}(\pi_1) = \operatorname{cl}(\pi_2)\) . We say that \(\operatorname{cl}(\pi)\) is the class of \(\pi\) .

Let c be the cardinal of A. Let \(\pi\) be a nonzero irreducible representation of A in a K-vector space E. Let \(\xi\) be a nonzero element of E. Since \(\pi(A)\xi = E\) (lemma 4), the dimension of E is \(\leqslant\mathfrak{c}\) (A, II, p.~97, corollary). The relation

“λ is a class of irreducible representations of A

in a K-vector space of dimension \(\leqslant\mathfrak{c}\gg\)

is collectivizing in \(\lambda\) (E, II, p.~3). Indeed, every vector space of dimension \(\leqslant \mathfrak{c}\) is isomorphic to a space \(\mathrm{K}^{\mathrm{B}}\) where B is a subset of A (A, II, p.~25, def. 10), and the assertion then follows from E, II, p.~47.

We denote by \(\widehat{\mathrm{A}}\) the set of classes of irreducible, nontrivial representations of A. According to what precedes, for every nontrivial irreducible representation \(\pi\) of A, there exists a unique representation \(\widehat{\pi} \in \widehat{\mathrm{A}}\) which is isomorphic to \(\pi\) .

The map from \(\hat{\mathrm{A}}\) into \(\mathrm{J(A)}\) which associates to \(\pi\) its kernel is surjective. If A is commutative, it follows from the fact that primitive ideals are regular maximal ideals that this map is a bijection.

We endow \(\hat{A}\) with the inverse-image topology under the map \(\hat{A} \rightarrow J(A)\) .

If A possesses a unit element, the spaces J(A) and \(\hat{\mathrm{A}}\) are quasi-compact.

It suffices to give the proof for J(A). Let \((\mathrm{T}_{j})\) be a family of closed subsets of J(A) whose intersection is empty. If the sum \(\sum_{j}\Upsilon(\mathrm{T}_{j})\) was different from A, then this sum would be contained in a maximal two-sided ideal I. The ideal I would be primitive (Lemma 7); since each \(T_{j}\) is closed, hence equal to \(\mathrm{V}(\Upsilon(\mathrm{T}_{j}))\) , we would have \(I\in T_{j}\) for all j, which contradicts the hypothesis. Thus we have \(\sum_{j}\Upsilon(\mathrm{T}_{j})=\mathrm{A}\) , and therefore we can write \(1=x_{1}+\cdots+x_{n}\) with \(n\geqslant1\) and \(x_{i}\in\Upsilon(\mathrm{T}_{j_{i}})\) for all i. This implies that \(\Upsilon(\mathrm{T}_{j_{1}})+\cdots+\Upsilon(\mathrm{T}_{j_{n}})=\mathrm{A}\) , whence \(T_{j_{1}}\cap\cdots\cap T_{j_{n}}=\varnothing\) .

Suppose the algebra A is commutative and unital. The Jacobson topology on J(A) is the topology induced on J(A) by the Zariski topology of the prime spectrum of A (AC, II, def. 4, p.~125).

Suppose that A is commutative and that K is a topological field. The canonical isomorphism from K onto \(\mathcal{L}(\mathrm{K})\) allows one to identify an element of \(\mathsf{X}(\mathsf{A})\) with a representation of A in the vector space K, which defines an injective map from \(\mathsf{X}(\mathsf{A})\) into \(\widehat{\mathsf{A}}\) . We can therefore identify \(\mathsf{X}(\mathsf{A})\) with a subset of \(\widehat{\mathsf{A}}\) .

Proposition 5. — The topology induced on \(\mathsf{X}(\mathsf{A})\) by that of \(\widehat{\mathsf{A}}\) is coarser than the topology of \(\mathsf{X}(\mathsf{A})\) .

Indeed, let T be a closed subset of \(\widehat{\mathbf{A}}\) . Then T is the set of \(\pi \in \widehat{\mathbf{A}}\) whose kernel contains a subset M of A. Therefore \(\mathrm{T} \cap \mathrm{X}(\mathrm{A})\) is the set of \(\chi \in \mathrm{X}(\mathrm{A})\) which vanish on M, that is, a closed subset of \(\mathrm{X}(\mathrm{A})\) . Whence the proposition.

In general, the topology of \(\mathsf{X}(\mathsf{A})\) does not coincide with the topology induced by the topology of \(\widehat{\mathsf{A}}\) (cf.~I, p.~193, exercise 6, c)).

\section{NORMED ALGEBRAS}

In this section, K denotes either of the fields R or C.

\subsection{Generalities}

Recall (cf.~TG, IX, p.~37, def. 9) that a normed algebra is an algebra A over K equipped with a norm \(x \mapsto \|x\|\) such that:

\[
\| x y \| \leqslant \| x \| \| y \|\tag{1}
\]

for all \(x, y \in A\) . If A is complete, A is said to be a Banach algebra.

Let A be a complete normable algebra over K. The topology of A can be defined by a norm satisfying (1) (cf.~TG, IX, p.~38). When A is equipped with the structure of a normed algebra defined by such a norm, we shall also say that the algebra A is a Banach algebra.

We also recall the following facts (cf.~TG, IX, p.~38--39):

\begin{enumerate}
\def\labelenumi{(\arabic{enumi})}
\item
  If A is a nonzero normed algebra and has a unit element e, then \(\|e\|\geqslant1\) .
\item
  Let A be a normed K-algebra. The opposite algebra of A, equipped with the same norm, is a normed algebra. The completion of A, equipped with the norm obtained by continuous extension of the norm of A, is a K-Banach algebra. Every subalgebra of A, equipped with the induced norm, is a normed algebra. If I is a closed two-sided ideal of A, the algebra A/I, equipped with the norm defined by \(\|\dot{x}\| = \inf_{x \in \dot{x}} \|x\|\) for all \(\dot{x} \in A/I\) is a normed algebra.
\item
  Let \((\mathrm{A}_{i})_{i\in\mathrm{I}}\) be a family of normed algebras and B the product algebra of the algebras \(A_{i}\) . The subalgebra of elements \((x_{i})_{i\in\mathrm{I}}\) of B such that \(\|(\boldsymbol{x}_{i})\|=\sup_{i\in\mathrm{I}}\|\boldsymbol{x}_{i}\|<+\infty\) is a normed algebra, called the product normed algebra of the algebras \((\mathrm{A}_{i})_{i\in\mathrm{I}}\) . If \(A_{i}\) is a Banach algebra for every i, then B is a Banach algebra.
\end{enumerate}

Let A be a normed algebra. Let \((y_{i})_{i\in\mathrm{I}}\) be a family of elements of A; the smallest closed subalgebra B of A containing the elements \(y_{i}\) is called the closed subalgebra of A generated by the \(y_{i}\) ; if B = A, one says that the \(y_{i}\) topologically generate the normed algebra A, or that the family \((y_{i})_{i\in\mathrm{I}}\) is a topological generating system of the normed algebra A.

Similarly, if A is a unital normed algebra, the smallest closed unital subalgebra containing the elements \(y_{i}\) is called the closed unital subalgebra of A generated by the \(y_{i}\) . If it is equal to A, one says that the \(y_{i}\) topologically generate the unital normed algebra A.

Let A be a normed K-algebra. Denote by \(\tilde{\mathbf{A}}\) the algebra obtained from A by adjoining an identity element. On \(\tilde{\mathbf{A}}\) , one defines a norm by setting \(\| (\lambda ,x)\| = |\lambda | + \| x\|\) for all \(\lambda \in \mathbf{K}\) and every \(x\in \mathbf{A}\) . We have

\[
\begin{array}{r l} & {\| (\lambda , x) (\mu , y) \| = | \lambda \mu | + \| x y + \mu x + \lambda y \|} \\ & {\qquad \leqslant | \lambda | | \mu | + \| x \| \| y \| + | \mu | \| x \| + | \lambda | \| y \|} \\ & {\qquad = \| (\lambda , x) \| \| (\mu , y) \|.} \end{array}
\]

Thus \(\widetilde{\mathsf{A}}\) becomes a normed algebra, called the unital normed algebra deduced from A by adjoining an identity element. The algebra A is identified with the closed two-sided ideal \(\{0\} \times \mathbf{A}\) of \(\widetilde{\mathsf{A}}\) .

DEFINITION 1. --- Let A be a normed algebra. For \(a \in \mathrm{A}\) , let \(\gamma_{a}\) and \(\delta_{a}\) be the maps \(x \mapsto ax\) , and \(x \mapsto xa\) from A to A. The map \(\gamma: a \mapsto \gamma_{a}\) is a representation of A in A, called the left regular representation of A. The map \(\delta: a \mapsto \delta_{a}\) is a representation of the opposite algebra of A in A, called the right regular representation of A.

Lemma 1. --- Let A be a normed algebra. The left regular representation of A and the right regular representation of A are continuous with norm \(\leqslant 1\) .

If the algebra A is unital, the maps \(x \mapsto \|\gamma_{x}\|\) and \(x \mapsto \|\delta_{x}\|\) are norms on A, defining the topology of A. They satisfy the inequality (1).

If the unital normed algebra A is nonzero, that is, if \(e \neq 0\) , we also have \(\|\gamma_{e}\| = \|\delta_{e}\| = 1\) .

It is immediate that \(\|\gamma_{x}\|\leqslant\|x\|\) and that \(\|\delta_{x}\|\leqslant\|x\|\) .

If A has an identity element \(e\) , then we have \(x = \gamma_x e = \delta_x e\) , hence:

\[
\| x \| \leqslant \| \boldsymbol {\gamma} _ {x} \| \cdot \| e \| \quad \| x \| \leqslant \| \boldsymbol {\delta} _ {x} \| \cdot \| e \|.\tag{2}
\]

In this case, \(x \mapsto \|\gamma_{x}\|\) and \(x \mapsto \|\delta_{x}\|\) are therefore norms equivalent to the norm of A. They satisfy (1) because \(\gamma\) and \(\delta\) are representations.

The last assertion follows from the fact that \(\gamma_{e} = \delta_{e} = Id_{A}\) .

\subsection{Examples}

\begin{enumerate}
\def\labelenumi{\arabic{enumi})}
\item
  Let E be a normed space. Equip E with the product defined by ab = 0 for all a and b in E; this defines on E a normed algebra structure.
\item
  Let X be a set. We denote by \(\mathcal{B}(X;K)\) the normed algebra of bounded functions on X with values in K, equipped with the norm

\[
\| f \| = \sup _ {x \in X} | f (x) |
\]

(TG, X, p.~22). It is a unital Banach algebra over K (TG, X, p.~21, cor. 1). It is commutative. Let \(f\) be an element of \(\mathcal{B}(\mathrm{X};\mathrm{K})\) . Then \(f\) is invertible in \(\mathcal{B}(\mathrm{X};\mathrm{K})\) if and only if one has

\[
\inf _ {x \in X} | f (x) | > 0.
\]

The spectrum of \(f\) is therefore the set of \(\lambda\in K\) such that

\[
\inf _ {x \in \mathrm{X}} | f (x) - \lambda | = 0,
\]

that is, the closure in K of the set \(f(\mathrm{X})\) of values of \(f\) .

\item
  Let X be a topological space. We denote by \(\mathcal{C}_{b}(\mathrm{X};\mathrm{K})\) the unital subalgebra of \(\mathcal{B}(\mathrm{X};\mathrm{K})\) of continuous and bounded functions on X with values in K, \(\mathcal{C}_{0}(\mathrm{X};\mathrm{K})\) the subalgebra of \(\mathcal{C}_{b}(\mathrm{X};\mathrm{K})\) consisting of functions that tend to 0 at infinity (cf.~INT, III, §1, n°2 and TG, X, p.~40, cor. 2). We recall that \(\mathcal{H}(\mathrm{X};\mathrm{K})\) denotes the subalgebra of \(\mathcal{C}_{b}(\mathrm{X};\mathrm{K})\) of continuous functions with compact support.
The algebras \(\mathcal{C}_{b}(\mathrm{X};\mathrm{K})\) and \(\mathcal{C}_{0}(\mathrm{X};\mathrm{K})\) are commutative Banach algebras over K; indeed, they are closed subspaces of \(\mathcal{B}(\mathrm{X};\mathrm{K})\) (TG, X, p.~21, cor. 2 and INT, III, §1, n°2).

Since the inverse of a continuous function that is nowhere zero is continuous, the subalgebra \(\mathcal{C}_{b}(\mathrm{X};\mathrm{K})\) is a full subalgebra of \(\mathcal{B}(\mathrm{X};\mathrm{K})\) . In particular, the spectrum of an element f of \(\mathcal{C}_{b}(\mathrm{X};\mathrm{K})\) is equal to \(\overline{f(\mathrm{X})}\) .

If X is discrete, one has \(\mathcal{C}_b(\mathrm{X};\mathrm{K}) = \mathcal{B}(\mathrm{X};\mathrm{K})\) .

If X is compact, one has \(\mathcal{C}_{b}(\mathrm{X};\mathrm{K}) = \mathcal{H}(\mathrm{X};\mathrm{K}) = \mathcal{C}(\mathrm{X};\mathrm{K})\) , the unital normed algebra \(\mathcal{C}(\mathrm{X};\mathrm{K})\) of continuous functions on X with values in K. In this case, an element \(f \in \mathcal{C}(\mathrm{X};\mathrm{K})\) is invertible if and only if f does not take the value 0, and the spectrum of f is equal to the set \(f(\mathrm{X})\) of values of f.

Suppose henceforth that X is not compact. The algebra \(\mathcal{C}_{0}(X;K)\) is then not unital; the algebra obtained from it by adjoining a unit element is identified with the subalgebra \(K \cdot 1 \oplus \mathcal{C}_{0}(X;K)\) of \(\mathcal{C}(X;K)\) formed by functions that have a limit at infinity. For an element of this subalgebra to be invertible, it is necessary and sufficient that it does not vanish and that its limit at infinity is not zero. It follows that for every \(f \in \mathcal{C}_{0}(X;K)\) , the spectrum of f is equal to \(f(X) \cup \{0\}\) .

The spectrum of \(f \in \mathcal{K}(\mathrm{X};\mathrm{K})\) is equal to \(f(\mathrm{X})\) (indeed, if \(\mathrm{X}\) is not compact, 0 belongs to \(f(\mathrm{X})\) ).

\item
  Let \(n \geqslant 0\) be an integer. Let \(A_{n}\) be the algebra of functions \(f: [0,1] \to K\) admitting continuous derivatives in \([0,1]\) up to order \(n\) , equipped with the norm
\[
\| f \| = \sum_ {k = 0} ^ {n} \frac {1}{k !} \sup _ {0 \leqslant t \leqslant 1} | f ^ {(k)} (t) |.
\]

If \(f,g\in \mathrm{A}_n\) , one has

\[
\begin{array}{l} \| f g \| = \sum_ {k = 0} ^ {n} \frac {1}{k !} \sup | (f g) ^ {(k)} (t) | = \sum_ {k = 0} ^ {n} \frac {1}{k !} \sup \Bigl | \sum_ {s = 0} ^ {k} \binom {k} {s} f ^ {(s)} (t) g ^ {(k - s)} (t) \Bigr | \\ \leqslant \sum_ {k = 0} ^ {n} \sum_ {s = 0} ^ {k} \frac {1}{s ! (k - s) !} \sup _ {0 \leqslant t \leqslant 1} | f ^ {(s)} (t) | \sup _ {0 \leqslant t \leqslant 1} | g ^ {(k - s)} (t) | = \| f \|   \| g \|, \end{array}
\]

by Leibniz's formula (FVR, I, p.~28, prop. 2), hence \(A_{n}\) is a unital commutative Banach algebra.

\item
  Let E be a Banach space whose norm we denote by p. The algebra \(\mathcal{L}(\mathrm{E})\) of continuous endomorphisms of E, equipped with the norm
\[
\| u \| = \sup _ {p (x) \leqslant 1} p (u (x))
\]

is a unital Banach algebra (EVT, III, p.~14 and p.~24, cor. 2; TG, X, p.~23, formula (3)).

\item
  Let G be a locally compact group. Denote by e its identity element. Let \(\mathcal{M}^{1}(G)\) be the Banach space of bounded complex measures on G (INT, III, p.~57). The convolution product (INT, VIII, p.~120, def. 1) equips \(\mathcal{M}^{1}(G)\) with a complex Banach algebra structure (INT, VIII, §3, n° 1, prop. 2) admitting as identity element the measure \(\varepsilon_{e}\) defined by the unit mass placed at the point e. If G is commutative, this Banach algebra is commutative. The space \(\mathcal{C}'(G)\) of measures with compact support is a subalgebra of \(\mathcal{M}^{1}(G)\) (loc. cit.).
\item
  Let G be a locally compact group equipped with a Haar measure \(\mu\) . Then \(\mathrm{L}_{\mathrm{K}}^{1}(\mathrm{G},\mu)\) is a Banach algebra for the convolution product (INT, VIII, prop. 12, p.~166). If K = C, the map defined by \(f \mapsto f \mu\) allows one to identify \(\mathrm{L}_{\mathbf{C}}^{1}(\mathrm{G},\mu)\) with a subalgebra of the Banach algebra \(\mathcal{M}^{1}(\mathrm{G})\) .
If G is commutative, the Banach algebra \(\mathrm{L}_{\mathrm{K}}^{1}(\mathrm{G},\mu)\) is commutative.

\item
  Take G = Z and K = C in Example 7. Then \(\mathrm{L}_{\mathbf{C}}^{1}(\mathbf{Z})\) is the complex commutative Banach algebra of sequences \((x_{n})_{n\in\mathbf{Z}}\) such that \(\sum_{n}|x_{n}|<+\infty\) , the product of elements \((x_{n})\) and \((y_{n})\) being \((z_{n})\) , where
\[
z _ {n} = \sum_ {k \in {\bf Z}} x _ {k} y _ {n - k}
\]

and the norm \(\|(x_{n})\|=\sum_{n}|x_{n}|\) . This algebra admits as its identity element the sequence \(\varepsilon=(\varepsilon_{n})\) such that \(\varepsilon_{0}=1\) and \(\varepsilon_{n}=0\) for \(n\neq0\) .

Let U denote the unit circle in C. If \(x = (c_{n})\) is an element of A, let \(\varphi(x)\) be the continuous function on U whose value at \(e^{it}\) is

\[
\varphi (x) (e ^ {i t}) = \sum_ {n \in {\bf Z}} c _ {n} e ^ {i n t}.
\]

One checks that \(\varphi\) is a morphism of \(\mathrm{L}_{\mathbf{C}}^{1}(\mathbf{Z})\) onto an algebra A of continuous functions on \(\mathbf{U}\) , multiplication in A being the usual multiplication. Integrating the equality term by term

\[
\left(\sum_ {m \in {\bf Z}} c _ {m} e ^ {i m t}\right) \cdot e ^ {- i n t} = \varphi (x) (e ^ {i t}) \cdot e ^ {- i n t},
\]

it comes

\[
c _ {n} = \frac {1}{2 \pi} \int_ {0} ^ {1} \varphi (x) (e ^ {i t}) e ^ {- i n t} d t.
\]

In particular, it follows that the morphism \(\varphi\) is injective. The algebra A, endowed with the norm induced by that of \(\mathrm{L}_{\mathbf{C}}^{1}(\mathbf{Z})\) via \(\varphi\) , is called the Banach algebra of absolutely convergent Fourier series. It admits as identity element the function \(1 = \varphi(\varepsilon)\) .

\item
  Let \(\Delta\) be the disk of complex numbers \(z\) satisfying \(|z| \leqslant 1\) . The algebra A of continuous functions on \(\Delta\) that are analytic in the interior of \(\Delta\) (VAR, R1, p.~26, 3.2.1) is equipped with the norm \(\|f\| = \sup_{z \in \Delta} |f(z)|\) . Then A is a commutative unital Banach algebra.
\end{enumerate}

\subsection{Spectral radius}

Lemma 2 (Fekete's Lemma). --- Let \((a_{n})_{n\geqslant1}\) be a sequence of real numbers. Suppose that

\[
a _ {n + m} \leqslant a _ {n} + a _ {m}
\]

for all \(n \geqslant 1\) and every \(m \geqslant 1\) . Then the sequence \((a_n / n)_{n \geqslant 1}\) converges and satisfies

\[
\lim _ {n \to + \infty} \frac {a _ {n}}{n} = \inf _ {n \geqslant 1} \frac {a _ {n}}{n}.
\]

Let us set \(a_0 = 0\) ; the inequality \(a_{n+m} \leqslant a_n + a_m\) remains valid for every \(n \geqslant 0\) and every \(m \geqslant 0\) . Fix an integer \(m \geqslant 1\) . For every integer \(n \geqslant 1\) , let \(q(n)\) and \(r(n)\) be the integers such that \(n = q(n)m + r(n)\) and \(0 \leqslant r(n) < m\) (E, III, p.~39, th. 1). The hypothesis then implies

\[
\frac {a _ {n}}{n} \leqslant \frac {a _ {q (n) m}}{n} + \frac {a _ {r (n)}}{n} \leqslant \frac {q (n) a _ {m}}{n} + \frac {a _ {r (n)}}{n} \leqslant \frac {q (n)}{n} a _ {m} + \frac {m}{n}.
\]

Letting \(n\) tend to \(+\infty\) , one deduces that \(\limsup_{n}(a_{n}/n) \leqslant a_{m}/m\) since \(q(n)/n \to 1/m\) . Since this holds for every \(m \geqslant 1\) , we therefore have

\[
\limsup _ {n \to + \infty} \frac {a _ {n}}{n} \leqslant \inf _ {m \geqslant 1} \frac {a _ {m}}{m} \leqslant \liminf _ {n \to + \infty} \frac {a _ {n}}{n}.
\]

These inequalities prove the convergence of the sequence \((a_{n}/n)_{n\geqslant1}\) as well as the formula \(\lim a_{n}/n = \inf_{n\geqslant1} a_{n}/n\) .

PROPOSITION 1. --- Let A be a normed algebra. For every \(x \in \mathrm{A}\) the sequence \((\|x^n\|^{1/n})_{n \geqslant 1}\) is convergent and its limit \(\varrho(x)\) is equal to \(\inf_{n \geqslant 1} \|x^n\|^{1/n}\) . Moreover, for every norm \(x \mapsto \|x\|_1\) defining the topology of A, one also has

\[
\varrho (x) = \lim _ {n \to + \infty} \| x ^ {n} \| _ {1} ^ {1 / n} = \inf _ {n \geqslant 1} \| x ^ {n} \| _ {1} ^ {1 / n}.
\]

If \(x\) is nilpotent, one has \(\| x^n\| _1^{1 / n} = 0\) for every integer \(n\) sufficiently large and every norm \(x\mapsto \| x\| _1\) defining the topology of A.

Suppose now that \(x\) is not nilpotent, and let us set \(\alpha_{n} = \| x^{n}\|\) . We have \(\alpha_{n} > 0\) for every integer \(n \geqslant 1\) , and \(\alpha_{n + m} \leqslant \alpha_{n}\alpha_{m}\) for all \(n, m \in \mathbf{N}\) by (1). Lemma 2, applied to the sequence \(a_n = \log(\alpha_n)\) shows the existence of the limit \(\varrho(x)\) and the formula \(\varrho(x) = \inf_{n > 0} \alpha_n^{1/n}\) .

Let \(x \mapsto \|x\|_1\) be a norm defining the topology of A. There exist real numbers \(a > 0\) and \(b > 0\) such that

\[
a \| x \| \leqslant \| x \| _ {1} \leqslant b \| x \|
\]

for all \(x \in A\) (EVT, II, p.~7, cor. 2). Consequently,

\[
a ^ {1 / n} \| x ^ {n} \| ^ {1 / n} \leqslant \| x ^ {n} \| _ {1} ^ {1 / n} \leqslant b ^ {1 / n} \| x ^ {n} \| ^ {1 / n},
\]

for all \(n \geqslant 1\) ; whence, passing to the limit, or taking the infimum, the equality

\[
\varrho (x) = \lim _ {n \to + \infty} \| x ^ {n} \| _ {1} ^ {1 / n} = \inf _ {n > 0} \| x ^ {n} \| _ {1} ^ {1 / n}.
\]

DEFINITION 2. --- For every element x of a normed algebra A, the real number

\[
\varrho (x) = \lim _ {n \to \infty} \| x ^ {n} \| ^ {1 / n} = \inf _ {n > 0} \| x ^ {n} \| ^ {1 / n}
\]

is called the spectral radius of x.

For every element \(x\) of A, we have

\[
\begin{array}{r l} (3) & \varrho (x) \leqslant \| x \| \\ (4) & \varrho (x ^ {n}) = \varrho (x) ^ {n}, \text { for every integer } n \geqslant 1. \end{array}
\]

DEFINITION 3. --- An element x of A is quasi-nilpotent if \(\varrho(x)=0\) .

This amounts to saying that, for any \(\lambda \in \mathbf{K}\) , the numbers \(\| (\lambda x)^n\|\) are bounded for \(n\geqslant 1\) , or equivalently that, for any \(\lambda \in \mathbf{K}\) , the sequence \((\lambda x)^n\) tends to 0 as \(n\to +\infty\) .

Remark 1. --- Let A be a normed algebra. If an element \(x \in \mathbf{A}\) satisfies \(\varrho(x) = \|x\|\) , one has \(\|x^n\| = \|x\|^n\) for all \(n \in \mathbf{N}\) , by (3) and (4).

Conversely, suppose that \(\|x^{2}\|=\|x\|^{2}\) for all \(x\in A\) . Then for every integer \(n\geqslant0\) , we have the equality \(\|x^{2^{n}}\|=\|x\|^{2^{n}}\) , hence \(\|x\|=\|x^{2^{n}}\|^{2^{-n}}\) ; as n tends to \(+\infty\) , we obtain \(\|x\|=\varrho(x)\) .

Remark 2. --- The function \(x \mapsto \varrho(x)\) on A, being the lower envelope of continuous functions \(x \mapsto \|x^n\|^{1/n}\) for \(n \geqslant 1\) , is upper semicontinuous (TG, IV, p.~31, cor.), but in general it is not continuous. It can even happen (cf.~exerc. 12 of I, p.~157) that a sequence of nilpotent elements of A tends to an element which is not quasi-nilpotent.

\subsection{Inverses}

Let A be a unital Banach algebra, whose identity element is denoted by 1. Recall (TG, IX, prop. 14, p.~40) that the group G of invertible elements of A is an open subset of A, that the topology induced on G by that of A is compatible with the group structure, and that the topological group G is complete.

PROPOSITION 2. --- Let A be a Banach algebra and x an element of A. The series \(\sum_{n=1}^{\infty}\lambda^{n}x^{n}\) considered as a power series in \(\lambda\) , has radius of convergence \(\varrho(x)^{-1}\) . If A is unital and if \(\varrho(x)<1\) , then 1-x is invertible and has inverse \(\sum_{n=0}^{\infty}x^{n}\) .

The series \(\sum_{n=1}^{\infty}\lambda^{n}x^{n}\) has radius of convergence

\[
(\limsup _ {n \to + \infty} \| x ^ {n} \| ^ {1 / n}) ^ {- 1} = \varrho (x) ^ {- 1}
\]

(cf.~VAR, R1, p.~23, 3.1.4). Suppose that A has a unit element. If \(\varrho(x) < 1\) , the series \(\sum_{n=0}^{\infty} x^n\) is therefore absolutely convergent. Since

\[
(1 - x) \left(\sum_ {n = 0} ^ {k} x ^ {n}\right) = \left(\sum_ {n = 0} ^ {k} x ^ {n}\right) (1 - x) = 1 - x ^ {k + 1}
\]

for every integer \(k \geqslant 0\) , the element 1 - x is invertible and its inverse is equal to \(\sum_{n=0}^{\infty} x^{n}\) .

COROLLARY 1. --- If A is a unital Banach algebra, then the group of invertible elements of A contains the open ball with center 1 and radius 1.

This is immediate since \(\|x\| < 1\) implies \(\varrho(x) < 1\) .

COROLLARY 2. --- Let A be a Banach algebra and I a maximal regular left (resp. right) ideal of A. Then I is closed.

Let ( \(\tilde{A}, e\) ) be the unital Banach algebra obtained from A by adjoining an identity element. There exists a maximal left (resp. right) ideal J of \(\tilde{A}\) such that \(J \cap A = I\) (A, VIII, p.~428, prop. 4). Then J is disjoint from the open ball with center e and radius 1 (cor. 1), and hence \(\overline{J} \neq \widetilde{A}\) . Since J is a maximal ideal, this implies that \(\overline{J} = J\) , and consequently \(I = J \cap A = \overline{J} \cap A\) is closed in A.

COROLLARY 3. --- The radical of a Banach algebra is closed.

Indeed, the radical is the intersection of the regular maximal left ideals (A, VIII, p.~430, def. 3).

PROPOSITION 3. --- Let A be a unital Banach algebra.

\begin{enumerate}
\def\labelenumi{\alph{enumi})}
\item
  If \(x \in A\) admits a left inverse (resp. right inverse) y, every element \(x' \in A\) such that \(\|x' - x\| < \|y\|^{-1}\) admits a left inverse (resp. right inverse).
\item
  The set of elements of A that are invertible (resp. left-invertible, resp. right-invertible) is open in A.
\item
  Let \((x_{n})\) be a sequence of elements of A admitting left inverses (resp. right inverses) \(y_{n}\) , such that \((x_{n})\) converges to an element \(x \in A\) and the sequence \((y_{n})\) is bounded. Then \(x\) is left-invertible (resp. right-invertible).
\end{enumerate}

It suffices to treat the case of left inverses; the case of right inverses follows by considering the opposite algebra.

Let \(x, y, x' \in \mathrm{A}\) such that \(yx = 1\) and \(\|x' - x\| < \|y\|^{-1}\) . We have

\[
\left\| 1 - y x ^ {\prime} \right\| = \left\| y x - y x ^ {\prime} \right\| \leqslant \left\| y \right\| \cdot \left\| x - x ^ {\prime} \right\| <   1,
\]

hence \(yx'\) is invertible: there exists \(z \in A\) such that \(z(yx') = 1\) . Thus the element \(x'\) is left-invertible with inverse zy. This proves assertion a), and assertion b) follows immediately.

  Let \((x_{n})\) be a sequence of elements of A admitting left inverses \(y_{n}\) such that \((x_{n})\) converges to an element \(x \in A\) and the sequence \((y_{n})\) is bounded. If \(M \geqslant 1\) is a real number such that \(\|y_{n}\| \leqslant M\) for all \(n \geqslant 1\) , then we have \(\|x_{n} - x\| < M^{-1} \leqslant \|y_{n}\|^{-1}\) for n large enough, and hence x admits a left inverse by a).

DEFINITION 4. --- Let A be a normed algebra, and let x be an element of A. Denote by \(\gamma_x\) and \(\delta_x\) the maps \(y \mapsto xy\) and \(y \mapsto yx\) from A to A. We say that \(x\) is a left (resp. right) topological zero divisor if \(\gamma_x\) (resp. \(\delta_x\) ) is not a homeomorphism from A onto \(\gamma_x(A)\) (resp. onto \(\delta_x(A)\) ).

Remark. --- By TG, IX, p.~36, cor. 2, x is a left (resp. right) topological zero divisor if and only if there exists a sequence \((z_{n})\) in A such that \(\|z_{n}\|=1\) and such that \(xz_{n}\) tends to 0 (resp. that \(z_{n}x\) tends to 0) when \(n\to+\infty\) .

A left (resp. right) topological zero divisor is a left (resp. right) zero divisor. Suppose that A is nonzero and unital. A left (resp. right) topological zero divisor x is not left (resp. right) invertible. Indeed, if for example yx = 1 and if \(xz_{n}\) tends toward 0, then \(z_{n} = y(xz_{n})\) tends to 0 and one cannot have \(\|z_{n}\| = 1\) for all n.

PROPOSITION 4. --- Let A be a unital Banach algebra. Let x be an element of A that is not left-invertible. If there exists a sequence \((x_{n})\) of left-invertible elements of A that converges to x, then x is a right topological zero divisor.

Let \(y_{n}\) be a left inverse of \(x_{n}\) . By prop. 3 (ii), \(\|y_{n}\|\) tends to \(+\infty\) . Let \(z_{n} = \|y_{n}\|^{-1}y_{n}\) . We have \(\|z_{n}\| = 1\) , and \(z_{n}x_{n} = \|y_{n}\|^{-1}\) tends to 0, hence \(z_{n}x = z_{n}x_{n} + z_{n}(x - x_{n})\) tends to 0. One then concludes with the aid of the remark following definition 4.

PROPOSITION 5. --- Let A be a unital Banach algebra and let B be a full subalgebra of A. Then \(\overline{\mathbf{B}}\) is a full subalgebra of A.

Indeed, let \(x\) be an element of \(\overline{\mathrm{B}}\) that is invertible in A, and let \((x_n)\) be a sequence of elements of B tending to \(x\) . Then, for \(n\) sufficiently large, \(x_n\) is invertible in A and \(x_n^{-1}\) tends to \(x^{-1}\) . Since the subalgebra B is full, we have \(x_n^{-1} \in \mathrm{B}\) , whence \(x^{-1} \in \overline{\mathrm{B}}\) .

Let A be a unital Banach algebra and let \((y_{i})_{i\in\mathrm{I}}\) be a family of elements of A. Let B be the full subalgebra of A generated by the elements \(y_{i}\) . Then \(\overline{B}\) is the smallest closed full subalgebra of A containing the \(y_{i}\) . It is called the closed full subalgebra generated by the elements \(y_{i}\) .

\subsection{Spectrum of an element in a normed algebra}

In this section, we assume that K = C.

THEOREM 1. --- Let A be a unital Banach algebra and let \(x \in \mathrm{A}\) .

\begin{enumerate}
\def\labelenumi{\alph{enumi})}
\item
  The set \(\mathrm{Sp}_{\mathrm{A}}(x)\) is a compact subset of C;
\item
  The spectral radius \(\varrho(x)\) is the radius of the smallest closed disk with center 0 in \(\mathbf{C}\) which contains \(\mathrm{Sp}_{\mathrm{A}}(x)\) ;
\item
  The resolvent \(\lambda \mapsto \mathrm{R}(x,\lambda) = (\lambda -x)^{-1}\) of \(x\) is holomorphic in \(\mathbf{C} - \mathrm{Sp}_{\mathrm{A}}(x)\) and zero at infinity. Moreover, for every integer \(k\geqslant 0\) , we have the
following formula:

\[
\frac {\partial^ {k}}{\partial \lambda^ {k}} \mathrm{R} (x, \lambda) = (- 1) ^ {k} k! \mathrm{R} (x, \lambda) ^ {k + 1};
\]

\item
  For every complex number \(\lambda\) such that \(|\lambda| > 1 / \varrho(x)\) , one has

\[
\mathrm{R} (x, \lambda) = \sum_ {n = 0} ^ {+ \infty} \lambda^ {- n - 1} x ^ {n}.
\]

\end{enumerate}

The complement of the spectrum \(\mathrm{Sp}_{\mathrm{A}}(x)\) is the inverse image of the group G of invertible elements of A under the continuous map \(\lambda \mapsto x - \lambda\) from C to A; by Proposition 3, b) of I, p.~23, the spectrum \(\mathrm{Sp}_{\mathrm{A}}(x)\) is a closed subset of C. Furthermore, let \(\lambda \in C\) such that \(|\lambda| > \varrho(x)\) . We have \(\lambda - x = \lambda(1 - \lambda^{-1}x)\) . Since \(\varrho(\lambda^{-1}x) = |\lambda|^{-1}\varrho(x) < 1\) , the element \(1 - \lambda^{-1}x\) , hence also the element \(\lambda - x\) , is invertible and

\[
\mathrm{R} (x, \lambda) = (\lambda - x) ^ {- 1} = \sum_ {n = 0} ^ {\infty} \lambda^ {- n - 1} x ^ {n}\tag{5}
\]

(I, p.~22, Prop. 2). In particular, \(\lambda \notin \mathrm{Sp}_{\mathrm{A}}(x)\) . This proves that \(\mathrm{Sp}_{\mathrm{A}}(x)\) is contained in the disk with center 0 and radius \(\varrho(x)\) . Consequently, \(\mathrm{Sp}_{\mathrm{A}}(x)\) is compact. This formula (5) also shows that the resolvent of \(x\) is defined and holomorphic in the complement of the closed disk \(\Delta_{\varrho(x)}\) of center 0 and radius \(\varrho(x)\) and tends to 0 at infinity.

Let \(\lambda_0\in \mathbf{C} - \mathrm{Sp}_{\mathrm{A}}(x)\) . Set \(y = \lambda_0 - x\) . Let \(\mu \in \mathbf{C}\) such that \(|\lambda_0 - \mu | < \| y^{-1}\|^{-1}\) . We have

\[
\mu - x = y - (\lambda_ {0} - \mu) = y (1 - (\lambda_ {0} - \mu) y ^ {- 1}),
\]

hence \(\mu - x\) is invertible and has as its inverse

\[
(\mu - x) ^ {- 1} = y ^ {- 1} \sum_ {n = 0} ^ {\infty} (\lambda_ {0} - \mu) ^ {n} y ^ {- n}\tag{6}
\]

By prop. 2 of I, p.~22. Hence the resolvent of \(x\) is defined and holomorphic in the open disk centered at \(\lambda_0\) and of radius \(\| y^{-1} \|^{-1}\) . Consequently, the resolvent of \(x\) is a holomorphic map from \(\mathbf{C} - \mathrm{Sp}_{\mathrm{A}}(x)\) into A.

Formula (1) of I, p.~4 implies \(\frac{\partial}{\partial\lambda}\mathrm{R}(x,\lambda) = -\mathrm{R}(x,\lambda)^2\) , whence, by induction on \(k\) ,

\[
\frac {\partial^ {k}}{\partial \lambda^ {k}} \mathrm{R} (x, \lambda) = (- 1) ^ {k} k! \mathrm{R} (x, \lambda) ^ {k + 1}.
\]

Let \(a > 0\) be a real number such that \(\mathrm{Sp}_{\mathrm{A}}(x)\) is contained in the closed disk \(\Delta_{a}\) of center 0 and radius \(a\) . The function \(\lambda \mapsto (\lambda^{-1} - x)^{-1}\) is then defined and holomorphic for \(0 < |\lambda| < a^{-1}\) and tends to 0 when \(\lambda\) tends to 0. The unique continuous function on the open disk centered at 0 with radius \(a^{-1}\) which extends this holomorphic function is then holomorphic (VAR, R1, 3.3.9), hence the radius of convergence of the series (5) defining it is \(\geqslant a^{-1}\) (VAR, R1, 3.2.9). By prop. 2 of I, p.~22, we therefore have \(a \geqslant \varrho(x)\) .

Remarks. --- 1) The spectrum of an element in a unital Banach algebra can be any arbitrary nonempty compact subset F of C (cf. Example 3 of I, p. 17; for A = C(F; C) and f ∈ A the canonical inclusion of F into C, one has \(\mathrm{Sp}_{\mathrm{A}}(f) = F\)).

\begin{enumerate}
\def\labelenumi{\arabic{enumi})}
\setcounter{enumi}{1}
\tightlist
\item
  Let A be a unital Banach algebra, and let \(x \in A\) . By Theorem 1 of I, p.~24, \(\mathbf{C} - \mathrm{Sp}_{\mathrm{A}}(x)\) is an open subset of C, hence is locally connected. Therefore the connected components of \(\mathbf{C} - \mathrm{Sp}_{\mathrm{A}}(x)\) are open. By Thm. 1, one of these connected components contains the set of \(\lambda \in C\) such that \(|\lambda| > \varrho(x)\) ; all the other connected components are therefore bounded.
\end{enumerate}

COROLLARY 1. --- Let A be a nonzero unital normed algebra. For every \(x \in \mathrm{A}\) the spectrum \(\mathrm{Sp}_{\mathrm{A}}(x)\) is nonempty.

Suppose first that A is complete. If one had \(\mathrm{Sp}(x)=\varnothing\) , the resolvent of x would be holomorphic in C and vanish at infinity, hence identically zero (VAR, R., 3.3.6, p.~29). Since \(\mathrm{R}(x,\lambda)=(\lambda-x)^{-1}\) is invertible, it would follow that 1=0 and hence that \(A=\{0\}\) .

In the general case, let \(\widehat{\mathrm{A}}\) be the completed algebra of A; the relation \(\mathrm{Sp}_{\mathrm{A}}(x) = \varnothing\) would entail \(\mathrm{Sp}_{\widehat{\mathrm{A}}} (x) = \varnothing\) , whence \(\widehat{\mathrm{A}} = \{0\}\) and \(\mathrm{A} = \{0\}\) .

COROLLARY 2 (Gelfand–Mazur theorem). --- Let A be a normed algebra over C. If A is a field, then A = C · 1.

If \(x \in A\) , there exists \(\lambda \in \mathbf{C}\) such that \(x - \lambda\) is non-invertible (cor. 1), hence \(x - \lambda = 0\) and \(x \in \mathbf{C} \cdot 1\) .

COROLLARY 3. --- Let A be a unital Banach algebra and let x be an invertible element of A such that \(\|x\|=\|x^{-1}\|=1\) . Then \(\mathrm{Sp}(x)\subset\mathbf{U}\) .

Let \(\Delta\) be the disk with center 0 and radius 1 in \(\mathbf{C}\) . By th. 1 b) and the fact that \(\varrho(x) \leqslant \|x\|\) , one has \(\operatorname{Sp}(x) \subset \Delta\) . Likewise, \(\operatorname{Sp}(x)^{-1} = \operatorname{Sp}(x^{-1}) \subset \Delta\) , whence the corollary (cf.~I, p.~2, remark 4).

COROLLARY 4. --- Let E be a complex Banach space, \(\mathcal{L}(\mathrm{E})\) the Banach algebra of continuous endomorphisms of E and A a nonzero subalgebra of \(\mathcal{L}(\mathrm{E})\) such that E is a simple A-pseudomodule (A, II, p.~176, Appendix).

\begin{enumerate}
\def\labelenumi{\alph{enumi})}
\item
  Let u be an endomorphism of E, not necessarily continuous, that commutes with A. Then u is a homothety.
\item
  Let \(u\) be an endomorphism of \(\mathbf{E}\) , not necessarily continuous. For every integer \(n \geqslant 1\) and for every \((\xi_1, \ldots, \xi_n) \in \mathrm{E}^n\) , there exists \(v \in A\) such that

\[
(v (\xi_ {1}), \dots , v (\xi_ {n})) = (u (\xi_ {1}), \dots , u (\xi_ {n})).
\]

\end{enumerate}

Let us show \(a\) ). Let \(\widetilde{\mathrm{A}}\) be the algebra obtained from A by adjoining an identity element. Since E is a simple A-pseudomodule, it is a simple \(\widetilde{\mathrm{A}}\)-module.

Let B be the commutant of \(\tilde{\mathsf{A}}\) in the ring of endomorphisms of the C-vector space E. The algebra B contains 1 and is the algebra of endomorphisms of the \(\tilde{\mathsf{A}}\)-module E. Since E is a simple \(\tilde{\mathsf{A}}\)-module, Schur's lemma (A, VIII, p.~43, corollary) shows that B is a field.

Let \(\xi_0\in \mathrm{E}\) be such that \(\mathrm{A}\xi_0\neq\{0\}\) . We therefore have \(\mathrm{A}\xi_0 = E\) . For every \(u\in \mathrm{B}\) , let \(\mathrm{A}_u\) be the set of \(v\in \mathrm{A}\) such that \(v(\xi_0) = u(\xi_0)\) . This set is nonempty, since \(\mathrm{A}\xi_0 = E\) . We then set

\[
\| u \| _ {\mathrm{B}} = \inf _ {v \in \mathrm{A} _ {u}} \| v \|.
\]

The map \(u \mapsto \|u\|_{\mathrm{B}}\) is a seminorm on B.

Let us show that this map is a norm. Let u be a nonzero element of B. For every \(v \in A_{u}\) , one has \(\|v\| \geqslant \|v(\xi_{0})\|/\|\xi_{0}\| = \|u(\xi_{0})\|/\|\xi_{0}\|\) , so that \(\|u\|_{B} \geqslant \|u(\xi_{0})\|/\|\xi_{0}\|\) . It therefore suffices to show that \(u(\xi_{0}) \neq 0\) . Let \(\xi_{1} \in E\) be such that \(u(\xi_{1}) \neq 0\) . Since \(A\xi_{0} = E\) , there exists \(w \in A\) such that \(\xi_{1} = w(\xi_{0})\) . Then, \(wu(\xi_{0}) = uw(\xi_{0}) = u(\xi_{1}) \neq 0\) , hence \(u(\xi_{0}) \neq 0\) .

On the other hand, let \(u\) and \(u'\) be elements of B. For every \(\varepsilon > 0\) , there exists \(v, v' \in A\) such that \(v(\xi_0) = u(\xi_0), v'(\xi_0) = u'(\xi_0)\) and \(\|v\| \leqslant \|u\|_{\mathrm{B}} + \varepsilon\) , \(\|v'\| \leqslant \|u'\|_{\mathrm{B}} + \varepsilon\) . Then we have \(vv'(\xi_0) = vu'(\xi_0) = u'v(\xi_0) = u'u(\xi_0)\) , whence

\[
\| u ^ {\prime} u \| _ {\mathrm{B}} \leqslant \| v ^ {\prime} v \| \leqslant \| v \| \| v ^ {\prime} \| \leqslant (\| u \| _ {\mathrm{B}} + \varepsilon) (\| u ^ {\prime} \| _ {\mathrm{B}} + \varepsilon),
\]

and finally \(\|u'u\|_{B} \leqslant \|u\|_{B} \|u'\|_{B}\) . This shows that B, equipped with the norm \(u \mapsto \|u\|_{B}\) , is a normed algebra. Since it is a field, Corollary 2 implies that \(B = C \cdot 1\) , which is the desired conclusion.

We prove \(b\) ). By \(a\) ), the commutant of A in \(\mathrm{End}_{\mathbf{C}}(\mathrm{E})\) is reduced to the homotheties of E. Its bicommutant is therefore \(\mathrm{End}_{\mathbf{C}}(\mathrm{E})\) . Assertion \(b\) ) therefore follows from the Jacobson density theorem (Theorem 1 of A, VIII, p.~434).

COROLLARY 5. --- Let A be a Banach algebra and \(x \in A\) .

\begin{enumerate}
\def\labelenumi{\alph{enumi})}
\item
  \(\mathrm{Sp}^{\prime}(x)\) is a compact subset of C;
\item
  The spectral radius \(\varrho(x)\) is the radius of the smallest closed disk centered at 0 in C that contains \(\mathrm{Sp}'(x)\) ;
\item
  For \(x\) to be quasi-nilpotent, it is necessary and sufficient that one has \(\mathrm{Sp}'(x) = \{0\}\) .
\end{enumerate}

The assertions \(a)\) and \(b)\) follow from th. 1 by considering the Banach algebra deduced from A by adjoining an identity element. The assertion \(c)\) follows from \(b)\) .

\subsection{Spectrum relative to a subalgebra}

In this section, we assume that K = C.

Lemma 3. --- Let \(X_{1}\) and \(X_{2}\) be compact subsets of C. If \(X_{2}\) is contained in \(X_{1}\) and if the boundary of \(X_{1}\) in C is contained in \(X_{2}\) , then \(X_{1}\) is the union of \(X_{2}\) and of certain bounded connected components of the complement of \(X_{2}\) in C.

Let U be a connected component of \(\mathbf{C}-X_{2}\) . Every boundary point of \(X_{1} \cap U\) in the open set U is also a boundary point of \(X_{1}\) in C, hence belongs to \(X_{2}\) by hypothesis; since \(U \cap X_{2} = \varnothing\) , we see that \(X_{1} \cap U\) has no boundary point in the space U. Since U is connected, the intersection \(X_{1} \cap U\) is either empty or equal to U (TG, I, p.~82, cor.), and the lemma follows.

PROPOSITION 6. — Let A be a unital Banach algebra and B a closed unital subalgebra of A. For every \(x \in B\) , one has \(\mathrm{Sp}_{\mathrm{B}}(x) \supset \mathrm{Sp}_{\mathrm{A}}(x)\) , and the boundary of \(\mathrm{Sp}_{\mathrm{A}}(x)\) in \(\mathbf{C}\) contains the boundary of \(\mathrm{Sp}_{\mathrm{B}}(x)\) in \(\mathbf{C}\) . In particular, if \(\mathrm{Sp}_{\mathrm{B}}(x) \subset \mathbf{R}\) , then we have \(\mathrm{Sp}_{\mathrm{B}}(x) = \mathrm{Sp}_{\mathrm{A}}(x)\) .

We have \(\mathrm{Sp}_{\mathrm{B}}(x) \supset \mathrm{Sp}_{\mathrm{A}}(x)\) (remark 6 of I, p.~3). If \(\lambda\) is a point on the boundary of \(\mathrm{Sp}_{\mathrm{B}}(x)\) in C, there exists a sequence \((\lambda_{n})\) of points exterior to \(\mathrm{Sp}_{\mathrm{B}}(x)\) tending to \(\lambda\) . Then \(x - \lambda_{n}\) is invertible in B and tends to \(x - \lambda\) , which is not invertible in B; hence \(x - \lambda\) is a left or right topological zero divisor in B (prop. 4 of I, p.~24), hence in A. Thus, \(\lambda \in \mathrm{Sp}_{\mathrm{A}}(x)\) . But since \(\mathrm{Sp}_{\mathrm{A}}(x) \subset \mathrm{Sp}_{\mathrm{B}}(x)\) , the complex number \(\lambda \in \mathrm{Fr}_{\mathbf{C}}(\mathrm{Sp}_{\mathrm{B}}(x))\) cannot be interior to \(\mathrm{Sp}_{\mathrm{A}}(x)\) , hence belongs to its boundary.

COROLLARY. --- The set \(\mathrm{Sp}_{\mathrm{B}}(x)\) is the union of \(\mathrm{Sp}_{\mathrm{A}}(x)\) and of certain bounded connected components of \(\mathbf{C}-\mathrm{Sp}_{\mathrm{A}}(x)\) .

This follows from prop. 6 and lemma 3.

This corollary will be completed by propositions 13 of I, p.~46 and 14 of I, p.~46.

\section{COMMUTATIVE BANACH ALGEBRAS}

In this section, the base field is C.

\subsection{Characters of a commutative Banach algebra}

THEOREM 1. --- Let A be a Banach algebra and let \(\chi\colon\mathbf{A}\to\mathbf{C}\) be an algebra homomorphism (cf.~I, p.~9). Then \(\chi\) is continuous, with norm at most 1. If A is unital and if \(\chi\) is a unital homomorphism, then \(\chi\) has norm 1.

Let us prove that \(|\chi(x)| \leqslant \|x\|\) for all \(x \in \mathbf{A}\) . Replacing A, if necessary, by the Banach algebra generated by \(x\) , we may assume that A is commutative; then \(\chi \in \mathbf{X}'(\mathbf{A})\) . For every \(x \in \mathbf{A}\) , one has \(\chi(x) \in \mathrm{Sp}_{\mathrm{A}}'(x)\) (I, p.~9, no.~7), hence \(|\chi(x)| \leqslant \varrho(x) \leqslant \|x\|\) (I, p.~28, cor. 5), whence the first assertion.

If, moreover, A is unital, the equality \(\chi(1)=1\) implies that \(\|\chi\|\geqslant1\) , whence the desired equality.

Remark. --- Let A be a commutative Banach algebra. It follows from this theorem that the topology of pointwise convergence on \(X'(A)\) coincides with the restriction to \(X'(A)\) of the weak topology \(\sigma(A', A)\) of the dual \(A'\) of A.

COROLLARY. --- Let A be a commutative Banach algebra. The space X'(A) is compact. The space X(A) is locally compact, and is compact if A has a unit element.

Let A' be the dual of the Banach space A. Its unit ball \(A_{1}^{\prime}\) is weakly compact (EVT, III, p.~17, cor. 2). By th. 1, we have \(\mathsf{X}^{\prime}(\mathsf{A}) \subset \mathsf{A}_{1}^{\prime}\) . Moreover, \(\mathsf{X}^{\prime}(\mathsf{A})\) is closed in A' for the weak topology, since it is the intersection of the weakly closed sets

\[
\mathrm{X} _ {x, y} = \left\{f \in \mathrm{A} ^ {\prime} \mid f (x y) - f (x) f (y) = 0 \right\} \quad (\text { for } x, y \in \mathrm{A}).
\]

It follows that \(X'(A)\) is compact and that \(X(A) = X'(A) - \{0\}\) is locally compact.

If A has a unit element 1, then \(\mathsf{X}(\mathsf{A})\) is the set of \(\chi \in \mathsf{X}'(\mathsf{A})\) such that \(\chi(1) = 1\) , and is therefore a closed subset, and hence compact, of \(\mathsf{X}'(\mathsf{A})\) .

Every compact space is homeomorphic to X(A) for a suitable commutative unital Banach algebra A (cf.~I, p.~32, cor. 2).

THEOREM 2. --- Let A be a commutative Banach algebra. The map \(\chi \mapsto \operatorname{Ker}(\chi)\) is a bijection from \(\mathsf{X}(\mathsf{A})\) onto the set \(\mathrm{J}(\mathrm{A})\) of the regular maximal ideals of A.

Let I be a regular maximal ideal of A. It is closed (I, p.~22, cor. 2), hence A/I is a Banach algebra. Since it is a field (A, VIII, p.~426, prop. 2), the Gelfand--Mazur theorem (I, p.~26, cor. 2) implies that A/I has dimension 1 over C. Hence I has codimension 1 in A. The theorem then follows from lemma 3 of I, p.~10.

It may happen that a nonzero commutative Banach algebra A, without unit element, has no maximal ideal; then \(X'(A)\) is reduced to \(\{0\}\) (cf.~I, p.~186, exercise 31).

This theorem shows that the sets \(J(A)\) and \(\hat{A}\) from I, p.~11, which are identified since A is commutative, can also be identified with the set \(X(A)\) . There is therefore reason to consider on \(X(A)\) the weak topology and the Jacobson topology, which is coarser (I, p.~15, prop. 5). The closed sets \(V(M)\) of the Jacobson topology, for \(M \subset A\) , are the sets

\[
\{\chi \in \mathrm{X} (\mathrm{A}) \mid \chi (\mathrm{M}) = 0 \}
\]

of \(\mathsf{X}(\mathsf{A})\) , which will sometimes still be denoted \(\mathsf{V}(\mathsf{M})\) . Similarly, for any subset \(\mathsf{M}\) of \(\mathsf{X}(\mathsf{A})\) , we shall denote by \(\Upsilon(\mathsf{M})\) the ideal intersection of the kernels of the \(\chi \in \mathsf{M}\) ; it is equal to the set \(\Upsilon(\mathsf{M}')\) defined in I, p.~13 for the subset \(\mathsf{M}' \subset \mathsf{J}(\mathsf{A})\) corresponding to M when one identifies \(\mathsf{J}(\mathsf{A})\) and \(\mathsf{X}(\mathsf{A})\) .

When a topological notion is used in X(A) without specifying which topology is meant, it will always be the weak topology.

\subsection{Continuous functions vanishing at infinity on a locally compact space}

In this section, X is a locally compact space. We denote by \(\mathcal{C}_{0}(X)\) the commutative Banach algebra of continuous complex-valued functions vanishing at infinity on X, equipped with the norm \(\|f\|=\sup_{x\in X}|f(x)|\) (example 3 of I, p.~17).

PROPOSITION 1. --- For every closed subset \(\Phi\) of X, let \(I_{\Phi}\) be the set of \(f \in \mathcal{C}_0(X)\) vanishing on \(\Phi\) . Then \(\Phi \mapsto I_{\Phi}\) is a bijection from the set of closed subsets of X onto the set of closed ideals of \(\mathcal{C}_0(X)\) .

The set \(\mathrm{I}_{\Phi}\) is a closed ideal of \(\mathcal{C}_0(\mathrm{X})\) .

Let \(\Phi \neq \Phi'\) be closed subsets of X. If necessary, by exchanging \(\Phi\) and \(\Phi'\) one may assume that there exists \(x \in \Phi'\) such that \(x \notin \Phi\) ; there then exists a function \(f \in \mathcal{C}_0(X)\) vanishing on \(\Phi\) and nonzero at \(x\) (TG, IX, p.~43, prop. 1). We have \(f \in I_{\Phi}\) and \(f \notin I_{\Phi'}\) , so that the map \(\Phi \mapsto I_{\Phi}\) is injective.

Let I be a closed ideal of \(\mathcal{C}_{0}(\mathrm{X})\) . Let \(\Phi\) be the set of \(x \in X\) such that \(f(x) = 0\) for all \(f \in I\) ; it is a closed subset of X, and we have \(I \subset I_{\Phi}\) . Let us prove that \(I_{\Phi} \subset I\) , which will imply that \(I = I_{\Phi}\) and will finish the proof of the proposition.

Let \(f \in \mathrm{I}_{\Phi}\) . For every real number \(\varepsilon > 0\) , let \(\mathrm{C}_{\varepsilon}\) be the set of \(x \in \mathrm{X}\) such that \(|f(x)| \geqslant \varepsilon\) . Since \(f\) tends to 0 at infinity, the set \(\mathrm{C}_{\varepsilon}\) is compact. Let \(x \in \mathrm{C}_{\varepsilon}\) ; since \(f(x) \neq 0\) and \(f \in \mathrm{I}_{\Phi}\) , one has \(x \notin \Phi\) ; by definition of \(\Phi\) , there then exists a function \(\varphi_x \in \mathrm{I}\) such that \(|\varphi_x(x)| > 1\) , hence such that \(|\varphi_x(y)| > 1\) for all \(y\) belonging to a neighborhood \(\mathrm{V}_x\) of \(x\) . The open sets \(\mathrm{V}_x \cap \mathrm{C}_{\varepsilon}\) cover \(\mathrm{C}_{\varepsilon}\) . Since the set \(\mathrm{C}_{\varepsilon}\) is compact, there exists a finite subset \(\mathrm{T}_{\varepsilon} \subset \mathrm{X}\) such that

\[
\mathrm{C} _ {\varepsilon} \subset \bigcup_ {x \in \mathrm{T} _ {\varepsilon}} \mathrm{V} _ {x}.
\]

Then the element

\[
g _ {\varepsilon} = \frac {1}{\varepsilon} \sum_ {x \in \mathrm{T} _ {\varepsilon}} \varphi_ {x} \overline {{\varphi_ {x}}} \geqslant 0
\]

of \(\mathcal{C}_0(\mathrm{X})\) belongs to I, and we have \(g_{\varepsilon} \geqslant \varepsilon^{-1}\) on \(\mathrm{C}_{\varepsilon}\) . The function

\[
f _ {\varepsilon} = \frac {f g _ {\varepsilon}}{1 + g _ {\varepsilon}}
\]

belongs to I. For \(x \notin \mathrm{C}_{\varepsilon}\) , one has

\[
\left| f (x) - f _ {\varepsilon} (x) \right| \leqslant 2 \varepsilon ,
\]

and for \(x\in \mathrm{C}_{\varepsilon}\) , one has

\[
| f (x) - f _ {\varepsilon} (x) | = \frac {| f (x) |}{1 + g _ {\varepsilon} (x)} \leqslant \varepsilon | f (x) |.
\]

Thus \(f_{\varepsilon}\) converges uniformly to f on X as \(\varepsilon\) tends to 0. We therefore have \(f \in \bar{I}\) , whence \(f \in I\) since I is closed.

COROLLARY 1. --- For every \(x \in X\) , let \(I_x\) be the set of \(f \in \mathcal{C}_0(X)\) vanishing at \(x\) . Then \(x \mapsto I_x\) is a bijection from \(X\) onto the set of maximal closed ideals of \(\mathcal{C}_0(X)\) . These ideals are regular.

This follows immediately from prop. 1.

Let X' denote the Alexandroff compactification of X, that is, the compact space obtained from X by adjoining a point at infinity \(\omega_{X}\) (TG, I, pp. 67–68). The algebra \(\mathcal{C}_{0}(X)\) is identified with the Banach algebra of continuous complex-valued functions on X' vanishing at \(\omega_{X}\) .

For every \(x \in X'\) , we denote by \(ev_{x}\) the character of \(\mathcal{C}_{0}(X)\) defined by \(\mathrm{ev}_{x}(f) = f(x)\) for all \(f \in \mathcal{C}_{0}(X)\) .

COROLLARY 2. --- The map \(x \mapsto \mathrm{ev}_x\) is a homeomorphism from \(\mathbf{X}'\) onto \(\mathsf{X}'(\mathcal{C}_0(\mathbf{X}))\) and its restriction to \(\mathbf{X}\) is a homeomorphism from \(\mathbf{X}\) onto \(\mathsf{X}(\mathcal{C}_0(\mathbf{X}))\) . Moreover, the weak topology and the Jacobson topology coincide on \(\mathsf{X}(\mathcal{C}_0(\mathbf{X}))\) .

The map ev: \(x \mapsto ev_{x}\) of \(X'\) into \(\mathsf{X}'(\mathcal{C}_{0}(\mathbf{X}))\) is injective. It is surjective by Cor. 1 and Th. 2 of I, p.~30. It is continuous, because for every function \(f \in \mathcal{C}_{0}(\mathbf{X})\) and every open set U of R, we have

\[
\mathrm{ev} ^ {- 1} (\{\chi \in \mathsf {X} ^ {\prime} (\mathcal {C} _ {0} (\mathsf {X})) \mid \chi (f) \in \mathrm{U} \}) = f ^ {- 1} (\mathrm{U})
\]

which is open in X. The map ev is therefore a homeomorphism since \(X'\) is compact. The restriction of ev to X is then a homeomorphism onto \(\mathsf{X}(\mathcal{C}_{0}(\mathsf{X}))\) .

If F is a weakly closed subset of \(\mathsf{X}(\mathcal{C}_{0}(\mathbf{X}))\) it corresponds, via the homeomorphism ev, to a closed subset \(\Phi\) of X; precisely, by Prop. 1, we have \(\mathrm{F} = \{\chi \in \mathsf{X}(\mathcal{C}_{0}(\mathbf{X})) \mid \mathrm{I}_{\Phi} \subset \mathrm{Ker} \, \chi\}\) which is closed for the Jacobson topology.

COROLLARY 3. --- Suppose X is compact. Then the map \(x \mapsto \mathrm{ev}_x\) is a homeomorphism from X onto \(\mathsf{X}(\mathcal{C}(\mathbb{X}))\) . The weak topology and the Jacobson topology coincide on \(\mathsf{X}(\mathcal{C}(\mathbb{X}))\) .
